% Lesson 22 — Grammar: Other forms of comparing, e.g. double, degree, equality, etc — Anglais Tle A — Cours signé OnBuch+
% Programme officiel MINESEC. Compiler avec : tectonic EN22-grammar-other-forms-of-comparing-e-g-double-degree-equa.tex
\newcommand{\DOCMATIERE}{Anglais}
\newcommand{\DOCNIVEAU}{Tle A}
\newcommand{\DOCMODULE}{Chapter 2 : Economic Life and Occupations}
\newcommand{\DOCLECON}{EN22}
\newcommand{\DOCTITRE}{Lesson 22 — Grammar: Other forms of comparing, e.g. double, degree, equality, etc}
% OnBuch+ — préambule « cours signés » ENRICHI (compilé avec tectonic / XeLaTeX)
% Système visuel « Pop » d'OnBuch+ : fond crème (jamais blanc), bordures encre
% épaisses, ombres « dures » décalées, titres Archivo Black en capitales.
% Version MATHÉMATIQUES : graphes (pgfplots/tikz), boîtes pédagogiques
% (définition, propriété, méthode, attention, le savais-tu, exercice résolu,
% exercice, TP, bilan), page de garde et signature OnBuch+ ancrée en bas.
%
% Macros attendues dans main.tex AVANT \input{preamble} :
%   \DOCMATIERE  \DOCNIVEAU  \DOCMODULE  \DOCLECON (numéro)  \DOCTITRE
\documentclass[11pt]{article}

\usepackage{fontspec}
\usepackage[english,french]{babel} % français = langue principale ; l'anglais via \eng{..} / textebox
\usepackage[a4paper,margin=2cm,top=3.4cm,bottom=2.8cm,headheight=1.4cm,footskip=1.3cm]{geometry}
\usepackage{amsmath,amssymb,amsfonts}
\usepackage{mathtools} % \xmapsto, \coloneqq... souvent utilisés par les modèles
\usepackage{cancel} % simplifications barrées (bilans de forces)
% \abs{z} (module, valeur absolue) et \norm{u} (norme), utilisés par les modèles
\providecommand{\abs}{}\let\abs\relax\DeclarePairedDelimiter{\abs}{\lvert}{\rvert}
\providecommand{\norm}{}\let\norm\relax\DeclarePairedDelimiter{\norm}{\lVert}{\rVert}
\usepackage[table]{xcolor} % [table] : fournit aussi \rowcolors (alternance de lignes)
\usepackage{pagecolor}
\usepackage{tikz}
\usetikzlibrary{arrows.meta,positioning,calc,shapes.geometric,shapes.misc,shapes.arrows,decorations.pathmorphing,decorations.pathreplacing,decorations.markings,patterns,shadows,mindmap,trees,fit,backgrounds,matrix,intersections,angles,quotes}
\usepackage{pgfplots}
\usepackage[version=4]{mhchem} % équations nucléaires \ce{^{235}_{92}U}
\usepackage[european,siunitx]{circuitikz} % schémas électriques
\pgfplotsset{compat=1.18}
\usepgfplotslibrary{fillbetween,groupplots} % aires entre courbes (calcul intégral)
\usepackage{enumitem}
\usepackage{fancyhdr}
% [most] charge toutes les bibliothèques tcolorbox (listings, minted, raster,
% documentation...) et gonfle inutilement la mémoire TeX sur un document long
% (~300 boîtes) : on ne charge que ce qui est réellement utilisé.
\usepackage[skins,breakable]{tcolorbox}
\usepackage{graphicx}
\usepackage{colortbl}
\usepackage{booktabs}
\usepackage{tabularx}
\usepackage{diagbox} % case d'en-tête diagonale des tableaux à double entrée
% Colonnes extensibles centrée / gauche / droite (C, L, R), souvent utilisées par les modèles
\newcolumntype{C}{>{\centering\arraybackslash}X}
\newcolumntype{L}{>{\raggedright\arraybackslash}X}
\newcolumntype{R}{>{\raggedleft\arraybackslash}X}
\usepackage{array}
\usepackage{multicol}
\usepackage{multirow}
\usepackage{siunitx}
% parse-numbers=false : accepte aussi \SI{2,5 \times 10^{-3}}{...}
\sisetup{locale=FR,output-decimal-marker={,},group-minimum-digits=4,parse-numbers=false}
\DeclareSIUnit\ohm{\ensuremath{\symup{\Omega}}} % U+2126 absent de la police
\DeclareSIUnit\division{div} % réglages d'oscilloscope (ms/div, V/div)
\DeclareSIUnit\tr{tr} % tours (vitesse de rotation, ex. \SI{1500}{\tr\per\minute})
\DeclareSIUnit\molar{M} % concentration molaire (ex. \SI{3}{\molar}, \SI{1}{\micro\molar})
\DeclareSIUnit\an{an} % année (le modèle confond parfois avec \year, un compteur TeX) — ex. \SI{5}{\kilo\meter\per\an}
\DeclareSIUnit\FCFA{FCFA} % franc CFA (ex. \SI{500}{\FCFA\per\kilo\gram})
\DeclareSIUnit\degreeBrix{\textdegree Brix} % teneur en sucre d'un sirop/jus (réfractométrie)
\DeclareSIUnit\month{mois} % le modèle utilise parfois \month, un compteur TeX, comme unité de durée
\usepackage{qrcode}
% \degree : commande fréquemment utilisée par les modèles pour un angle ou une
% température en dehors de \SI{}{} (non fournie par les paquets chargés).
\providecommand{\degree}{^{\circ}}
% \qed (fin de démonstration), utilisé par les modèles sans amsthm
\providecommand{\qed}{\hfill$\square$}
% \barre : double barre des valeurs interdites dans un tableau de variations (tkz-tab)
\providecommand{\barre}{\ensuremath{\|}}
% environnement proof (amsthm non chargé) utilisé par les modèles
\ifdefined\proof\else\newenvironment{proof}[1][Démonstration]{\par\noindent\textit{#1.}\ }{\qed\par}\fi
\usepackage{lastpage}
\usepackage{hyperref}
\hypersetup{hidelinks}

% ============================================================
% Palette « Pop » OnBuch+ — mêmes hex que la classe P (plus_ui.dart)
% ============================================================
\definecolor{poporange}{HTML}{F59321}
\definecolor{popdark}{HTML}{E07A0C}
\definecolor{popcream}{HTML}{FFF6E8}
\definecolor{popink}{HTML}{1C1714}
\definecolor{poppurple}{HTML}{7A5AE0}
\definecolor{popgreen}{HTML}{2E9E6A}
\definecolor{poppink}{HTML}{E0457B}
\definecolor{popblue}{HTML}{2F7DE1}
\definecolor{popgold}{HTML}{C98A12}
\definecolor{popmuted}{HTML}{8A7F76}
\definecolor{popline}{HTML}{EADFD2}
\definecolor{popgray}{HTML}{8A7F76}
\colorlet{popgrey}{popgray}
% Alias tolérants (le modèle nomme parfois les couleurs différemment)
\colorlet{popwhite}{white}
\colorlet{popblack}{popink}
\colorlet{popred}{poppink}
\colorlet{popyellow}{popgold}
% Teintes claires (fonds de tableaux/encadrés) — le modèle nomme parfois
% les couleurs avec un suffixe L (ex. popblueL) sans les avoir définies.
\colorlet{poporangeL}{poporange!20}
\colorlet{popdarkL}{popdark!20}
\colorlet{poppurpleL}{poppurple!20}
\colorlet{popgreenL}{popgreen!20}
\colorlet{poppinkL}{poppink!20}
\colorlet{popblueL}{popblue!20}
\colorlet{popgoldL}{popgold!20}
\colorlet{popredL}{popred!20}
\colorlet{popgrayL}{popgray!20}
% Teintes claires pour fonds de tableaux / zones
\colorlet{popblueL}{popblue!10!white}
\colorlet{poporangeL}{poporange!14!white}
\colorlet{popgreenL}{popgreen!12!white}
\colorlet{poppurpleL}{poppurple!12!white}
\colorlet{poppinkL}{poppink!12!white}
\colorlet{popgoldL}{popgold!14!white}

\pagecolor{popcream}

% ============================================================
% Typographie
% ============================================================
\defaultfontfeatures{Ligatures=TeX}
\setmainfont{PlusJakartaSans-Medium.ttf}[Path=../../../fonts/,BoldFont=PlusJakartaSans-Bold.ttf]
% Symboles, API et emojis absents de la police principale : repli sur DejaVu Sans (caractères actifs)
\newfontface\symfont{DejaVuSans.ttf}[Path=../../../fonts/]
\makeatletter
\newcommand\symactive[1]{\catcode#1=13 \begingroup\lccode`\~=#1 \lowercase{\endgroup\def~}{{\symfont\char#1}}}
\newcommand\symblank[1]{\catcode#1=13 \begingroup\lccode`\~=#1 \lowercase{\endgroup\def~}{}}
\newcommand\symhyph[1]{\catcode#1=13 \begingroup\lccode`\~=#1 \lowercase{\endgroup\def~}{\mbox{-}}}
\newcount\symc
\newcommand\symrange[2]{\symc=#1 \@whilenum\symc<#2 \do{\expandafter\symactive\expandafter{\the\symc}\advance\symc by 1 }}
\newcommand\symblankrange[2]{\symc=#1 \@whilenum\symc<#2 \do{\expandafter\symblank\expandafter{\the\symc}\advance\symc by 1 }}
\makeatother
\symrange{"0250}{"02FF}\symactive{"2713}\symactive{"2714}\symactive{"2717}\symactive{"2718}\symactive{"2610}\symactive{"2611}\symactive{"2612}
\symhyph{"2011}\symhyph{"2E1A}\symblankrange{"1F300}{"1FAFF}\symblank{"FE0F}
% Maths en sans-serif (Fira Math) avec chiffres et lettres droites de Plus
% Jakarta, pour que les formules se fondent dans le texte.
\usepackage[math-style=ISO,bold-style=ISO]{unicode-math}
\setmathfont{FiraMath-Regular.otf}[Path=../../../fonts/]
\setmathfont{PlusJakartaSans-Medium.ttf}[Path=../../../fonts/,range={up/num,up/{Latin,latin},\mathup}]
\usepackage{icomma} % virgule décimale sans espace en mode math
% Caractères mathématiques Unicode absents des polices de texte (Plus Jakarta,
% Archivo Black) : rendus par leur équivalent mathématique, en texte comme en
% formule — sinon ils s'affichent en carré vide (ex. « Arithmétique dans ℤ »).
\usepackage{newunicodechar}
\newunicodechar{ℝ}{\ensuremath{\mathbb{R}}}
\newunicodechar{ℤ}{\ensuremath{\mathbb{Z}}}
\newunicodechar{ℂ}{\ensuremath{\mathbb{C}}}
\newunicodechar{ℕ}{\ensuremath{\mathbb{N}}}
\newunicodechar{ℚ}{\ensuremath{\mathbb{Q}}}
\newunicodechar{𝔻}{\ensuremath{\mathbb{D}}}
\newunicodechar{α}{\ensuremath{\alpha}}
\newunicodechar{β}{\ensuremath{\beta}}
\newunicodechar{γ}{\ensuremath{\gamma}}
\newunicodechar{δ}{\ensuremath{\delta}}
\newunicodechar{ε}{\ensuremath{\varepsilon}}
\newunicodechar{θ}{\ensuremath{\theta}}
\newunicodechar{λ}{\ensuremath{\lambda}}
\newunicodechar{μ}{\ensuremath{\mu}}
\newunicodechar{σ}{\ensuremath{\sigma}}
\newunicodechar{τ}{\ensuremath{\tau}}
\newunicodechar{φ}{\ensuremath{\varphi}}
\newunicodechar{ψ}{\ensuremath{\psi}}
\newunicodechar{ω}{\ensuremath{\omega}}
\newunicodechar{Σ}{\ensuremath{\Sigma}}
% majuscules produites par \MakeUppercase dans les titres (α-aminés -> Α)
\newunicodechar{Α}{\ensuremath{\alpha}}
\newunicodechar{Β}{\ensuremath{\beta}}
\newunicodechar{Γ}{\ensuremath{\Gamma}}
\newunicodechar{Δ}{\ensuremath{\Delta}}
\newunicodechar{Ε}{\ensuremath{\varepsilon}}
\newunicodechar{Θ}{\ensuremath{\Theta}}
\newunicodechar{Λ}{\ensuremath{\Lambda}}
\newunicodechar{Μ}{\ensuremath{\mu}}
\newunicodechar{Τ}{\ensuremath{\tau}}
\newunicodechar{Φ}{\ensuremath{\Phi}}
\newunicodechar{Ψ}{\ensuremath{\Psi}}
\newunicodechar{Ω}{\ensuremath{\Omega}}
\newunicodechar{↦}{\ensuremath{\mapsto}}
\newunicodechar{∈}{\ensuremath{\in}}
\newunicodechar{∉}{\ensuremath{\notin}}
\newunicodechar{∧}{\ensuremath{\wedge}}
\newunicodechar{⇔}{\ensuremath{\Leftrightarrow}}
\newunicodechar{⟺}{\ensuremath{\Longleftrightarrow}}
\newunicodechar{⇒}{\ensuremath{\Rightarrow}}
\newunicodechar{⟹}{\ensuremath{\Longrightarrow}}
\newunicodechar{∘}{\ensuremath{\circ}}
\newunicodechar{∪}{\ensuremath{\cup}}
\newunicodechar{∩}{\ensuremath{\cap}}
\newunicodechar{≡}{\ensuremath{\equiv}}
\newunicodechar{⊂}{\ensuremath{\subset}}
\newunicodechar{⊥}{\ensuremath{\perp}}
\newunicodechar{∃}{\ensuremath{\exists}}
\newunicodechar{∀}{\ensuremath{\forall}}
\newunicodechar{∛}{\ensuremath{\sqrt[3]{\,}}}
\newunicodechar{∜}{\ensuremath{\sqrt[4]{\,}}}
\newunicodechar{⃗}{\ensuremath{{}^{\to}}}
\newunicodechar{ⁿ}{\ifmmode^{n}\else\textsuperscript{\ensuremath{n}}\fi}
\newunicodechar{ˣ}{\ifmmode^{x}\else\textsuperscript{\ensuremath{x}}\fi}
\newunicodechar{ᵇ}{\ifmmode^{b}\else\textsuperscript{\ensuremath{b}}\fi}
\newunicodechar{ʳ}{\ifmmode^{r}\else\textsuperscript{\ensuremath{r}}\fi}
\newunicodechar{ᵘ}{\ifmmode^{u}\else\textsuperscript{\ensuremath{u}}\fi}
\newunicodechar{ʸ}{\ifmmode^{y}\else\textsuperscript{\ensuremath{y}}\fi}
\newunicodechar{ᵃ}{\ifmmode^{a}\else\textsuperscript{\ensuremath{a}}\fi}
\newunicodechar{ᵉ}{\ifmmode^{e}\else\textsuperscript{\ensuremath{e}}\fi}
\newunicodechar{ᵏ}{\ifmmode^{k}\else\textsuperscript{\ensuremath{k}}\fi}
\newunicodechar{ᵖ}{\ifmmode^{p}\else\textsuperscript{\ensuremath{p}}\fi}
\newunicodechar{ᴺ}{\ifmmode^{N}\else\textsuperscript{\ensuremath{N}}\fi}
\newunicodechar{ᶜ}{\ifmmode^{c}\else\textsuperscript{\ensuremath{c}}\fi}
\newunicodechar{ʰ}{\ifmmode^{h}\else\textsuperscript{\ensuremath{h}}\fi}
\newunicodechar{ᵠ}{\ifmmode^{q}\else\textsuperscript{\ensuremath{q}}\fi}
\newunicodechar{ₙ}{\ifmmode_{n}\else\textsubscript{\ensuremath{n}}\fi}
\newunicodechar{ₐ}{\ifmmode_{a}\else\textsubscript{\ensuremath{a}}\fi}
\newunicodechar{ᵢ}{\ifmmode_{i}\else\textsubscript{\ensuremath{i}}\fi}
\newunicodechar{ᵣ}{\ifmmode_{r}\else\textsubscript{\ensuremath{r}}\fi}
\newunicodechar{ₖ}{\ifmmode_{k}\else\textsubscript{\ensuremath{k}}\fi}
\newunicodechar{ᵦ}{\ifmmode_{\beta}\else\textsubscript{\ensuremath{\beta}}\fi}
\newunicodechar{ₓ}{\ifmmode_{x}\else\textsubscript{\ensuremath{x}}\fi}
\newfontfamily\popdisplay{ArchivoBlack-Regular.ttf}[Path=../../../fonts/]
\newfontfamily\poplogo{SpaceGrotesk-Bold.ttf}[Path=../../../fonts/]
\newfontfamily\popsemi{PlusJakartaSans-SemiBold.ttf}[Path=../../../fonts/]
\newfontfamily\popxbold{PlusJakartaSans-ExtraBold.ttf}[Path=../../../fonts/]
\newfontfamily\popxboldls{PlusJakartaSans-ExtraBold.ttf}[Path=../../../fonts/,LetterSpace=3.0]

\setlist[enumerate]{leftmargin=1.7em,itemsep=5pt,topsep=4pt}
\setlist[itemize]{leftmargin=1.7em,itemsep=4pt,topsep=4pt,label={\color{poporange}\textbullet}}
\setlength{\parindent}{0pt}
\setlength{\parskip}{5pt}
\renewcommand{\arraystretch}{1.25}

% pgfplots : style Pop par défaut
\pgfplotsset{
  every axis/.append style={axis line style={popink,line width=1pt},tick style={popink},
    grid style={popline,line width=0.5pt},label style={font=\small},tick label style={font=\footnotesize}},
  popcurve/.style={poporange,line width=1.8pt,no markers,smooth},
}
% Même style disponible dans un \draw TikZ simple (hors pgfplots)
\tikzset{popcurve/.style={poporange,line width=1.8pt}}
\tikzset{popgrid/.style={popline,very thin}}
\colorlet{popcurve}{poporange} % le nom du style est parfois utilisé comme couleur

% ============================================================
% Pill / squiggle
% ============================================================
\newcommand{\poppill}[3]{%
  \tikz[baseline=-0.55ex]\node[draw=popink,line width=1.2pt,rounded corners=2.6pt,%
    fill=#2,inner xsep=6.5pt,inner ysep=3pt]{%
    \popxboldls\fontsize{7}{7}\selectfont\color{#3}\MakeUppercase{#1}};%
}
\newcommand{\popsquiggle}[1]{%
  \tikz[baseline=-0.4ex]{
    \draw[#1,line width=2.6pt,line cap=round]
      plot [smooth] coordinates {(0,0.09) (0.34,0.01) (0.68,0.21) (1.02,0.01) (1.36,0.10)};
  }%
}
% Mot-clé surligné (marqueur orange sous le texte)
\newcommand{\cle}[1]{{\popxbold\color{popdark}#1}}

% ============================================================
% En-tête / pied
% ============================================================
\pagestyle{fancy}
\fancyhf{}
\renewcommand{\headrulewidth}{0pt}
\renewcommand{\footrulewidth}{0pt}
\lhead{\raisebox{-0.15ex}{\poplogo\fontsize{13}{13}\selectfont\color{popink}OnBuch\color{poporange}+}}
\rhead{\poppill{\DOCMATIERE\ \textbullet\ \DOCNIVEAU\ \textbullet\ Leçon \DOCLECON}{white}{popink}}
\lfoot{\popxbold\fontsize{7.6}{7.6}\selectfont\color{popmuted}\MakeUppercase{Cours signé OnBuch+}\ \textbullet\ {\popsemi\DOCTITRE}}
\rfoot{\tikz[baseline=-0.6ex]\node[rounded corners=8.5pt,fill=popink,minimum height=17pt,minimum width=17pt,inner xsep=5pt,inner sep=0]{\popxbold\fontsize{8}{8}\selectfont\color{popcream}\thepage\,/\,\pageref*{LastPage}};}

% ============================================================
% Titres
% ============================================================
\newcommand{\chaptitre}[2]{%
  \begin{tcolorbox}[enhanced,colback=poporange,colframe=popink,boxrule=2.6pt,arc=11pt,
      drop shadow=popink,left=15pt,right=15pt,top=12pt,bottom=11pt,before skip=3pt,after skip=18pt]
    \poppill{#2}{popink}{popcream}\par\vspace{9pt}
    {\raggedright\hyphenpenalty=10000\exhyphenpenalty=10000\popdisplay\fontsize{22}{24}\selectfont\color{popcream}\MakeUppercase{#1}\par}
  \end{tcolorbox}
}
% Section numérotée : pastille orange avec le numéro + titre Archivo
\newcounter{popsec}
\newcommand{\coursec}[1]{%
  \stepcounter{popsec}\setcounter{popsub}{0}%
  \par\vspace{16pt}\noindent
  \begin{minipage}{\linewidth}
  \tikz[baseline=-0.7ex]\node[draw=popink,line width=1.6pt,fill=poporange,rounded corners=4pt,minimum size=22pt,inner sep=0]{\popdisplay\fontsize{12}{12}\selectfont\color{popcream}\thepopsec};%
  \hspace{8pt}{\popdisplay\fontsize{15}{17}\selectfont\color{popink}\MakeUppercase{#1}}\par
  \vspace{4pt}\noindent{\color{poporange}\rule{36pt}{3.6pt}}
  \end{minipage}\par\nopagebreak\vspace{8pt}
}
\newcounter{popsub}
\newcommand{\courssub}[1]{%
  \stepcounter{popsub}%
  \par\vspace{9pt}\noindent{\popxbold\fontsize{11.5}{13}\selectfont\color{popink}{\color{poporange}\thepopsec.\thepopsub}\hspace{6pt}#1}\par\nopagebreak\vspace{4pt}
}

% ============================================================
% Boîtes pédagogiques — même recette : fond blanc, bordure épaisse colorée,
% ombre dure encre, badge plein. Titre optionnel [#1].
% ============================================================
\tcbset{popbase/.style={breakable,enhanced,colback=white,boxrule=2.3pt,arc=9pt,
  shadow={3pt}{-3pt}{0pt}{popink},left=14pt,right=14pt,top=11pt,bottom=11pt,before skip=11pt,after skip=15pt,
  fontupper=\popsemi\fontsize{10.6}{14.5}\selectfont\color{popink},
  fontlower=\popsemi\fontsize{10.6}{14.5}\selectfont\color{popink}}}
% Coche dessinée (le glyphe ✓ n'existe pas dans les polices du document)
\newcommand{\popcheck}[1]{\tikz[baseline=-0.5ex]\draw[#1,line width=1.6pt,line cap=round,line join=round](-0.13,0.0)--(-0.04,-0.09)--(0.13,0.1);}
\providecommand{\coche}{\popcheck{popgreen}}
\providecommand{\resultat}[1]{\par\smallskip\noindent\fcolorbox{popgreen}{popgreenL}{\parbox{\dimexpr\linewidth-2\fboxsep-2\fboxrule\relax}{\textbf{Résultat.} #1}}\par\smallskip}
\newcommand{\popboxhead}[3]{\poppill{#1}{#2}{white}\ifx\relax#3\relax\else\hspace{6pt}{\popxbold\fontsize{10.6}{12}\selectfont\color{popink}#3}\fi\par\vspace{7pt}}

\newtcolorbox{aretenir}[1][]{popbase,colframe=popgreen,before upper={\popboxhead{À retenir}{popgreen}{#1}}}
% Texte en anglais (document de lecture, dialogue, extrait) : cadre
% d'étude. Un titre optionnel (source, auteur) s'affiche dans l'en-tête.
\newtcolorbox{textebox}[1][]{popbase,colframe=popblue,before upper={\popboxhead{Text}{popblue}{#1}\begin{otherlanguage}{english}},after upper={\end{otherlanguage}}}
% Marques ✓ / ✗ (forme correcte / fautive) souvent utilisées par les modèles
\providecommand{\cross}{\textcolor{poppink}{\ensuremath{\times}}}
\providecommand{\xmark}{\cross}\providecommand{\crossmark}{\cross}\providecommand{\cmark}{\ensuremath{\checkmark}}
% Ligne de réponse / justification à compléter (inventée par les modèles)
\providecommand{\justification}[1]{\par\noindent\textit{Justification~:} \dotfill\par}
% \textipa (package tipa non chargé) : la police ne couvre pas l'API -> simple passage
\providecommand{\textipa}[1]{#1}
% Soulignement (ulem non chargé) : \uline, \ul
\providecommand{\uline}[1]{\underline{#1}}\providecommand{\ul}[1]{\underline{#1}}
% \glossaire{\item ...} inventée par les modèles : liste de glossaire
\providecommand{\glossaire}[1]{\begin{itemize}#1\end{itemize}}
% \alert (beamer) inventée par les modèles : texte mis en évidence
\providecommand{\alert}[1]{\textbf{\textcolor{poppink}{#1}}}
% variantes ASCII d'ordinaux écrites en macro
\providecommand{\iere}{\textsuperscript{re}}\providecommand{\ieme}{\textsuperscript{e}}\providecommand{\ere}{\textsuperscript{re}}\providecommand{\eme}{\textsuperscript{e}}\providecommand{\er}{\textsuperscript{er}}\providecommand{\ie}{\textsuperscript{e}}
% Ordinaux français écrits en macro par les modèles : 1\ière, 3\ième
\newcommand{\ière}{\textsuperscript{re}}\newcommand{\ième}{\textsuperscript{e}}
% \exemple (inventée par les modèles) : étiquette « Exemple : »
\providecommand{\exemple}{\textbf{Exemple~:}\ }
% \square utilisable aussi hors mode math (cases à cocher)
\AtBeginDocument{\let\squareMath\square\renewcommand{\square}{\ensuremath{\squareMath}}}
% Mot ou phrase en anglais à l'intérieur d'un paragraphe français
\newcommand{\eng}[1]{\ifmmode\text{\textit{#1}}\else\begingroup\selectlanguage{english}\itshape #1\endgroup\fi}
\let\esp\eng % alias toléré
% flèches d'intonation inventées par les modèles
\providecommand{\textsearrow}{}\renewcommand{\textsearrow}{\ensuremath{\searrow}}\providecommand{\textnearrow}{}\renewcommand{\textnearrow}{\ensuremath{\nearrow}}
% \fr{..} : mot français (inventé par les modèles)
\providecommand{\fr}[1]{\textit{#1}}
\newtcolorbox{exemplebox}[1][]{popbase,colframe=poporange,before upper={\popboxhead{Exemple}{poporange}{#1}}}
\newtcolorbox{definition}[1][]{popbase,colframe=popblue,before upper={\popboxhead{Définition}{popblue}{#1}}}
\newtcolorbox{propriete}[1][]{popbase,colframe=poppurple,before upper={\popboxhead{Propriété}{poppurple}{#1}}}
\newtcolorbox{methode}[1][]{popbase,colframe=popgold,before upper={\popboxhead{Méthode}{popgold}{#1}}}
\newtcolorbox{attention}[1][]{popbase,colframe=poppink,before upper={\popboxhead{Attention}{poppink}{#1}}}
\newtcolorbox{experience}[1][]{popbase,colframe=popblue,colback=popblueL,before upper={\popboxhead{Expérience}{popblue}{#1}}}
% « Le savais-tu ? » : carte violette pleine (contexte, culture, Cameroun)
\newtcolorbox{savaistu}[1][]{popbase,colback=poppurple,colframe=popink,
  fontupper=\popsemi\fontsize{10.4}{14.2}\selectfont\color{white},
  before upper={\poppill{Le savais-tu\,?}{white}{poppurple}\ifx\relax#1\relax\else\hspace{6pt}{\popxbold\color{white}#1}\fi\par\vspace{7pt}}}
% Exercice résolu : énoncé en haut, \tcblower puis la solution sur fond crème
\newcounter{popexo}\newcounter{popexr}
\newtcolorbox{exoresolu}[1][]{popbase,colframe=poporange,colbacklower=popcream,
  segmentation style={popink,line width=1.2pt,dashed},
  before upper={\stepcounter{popexr}\popboxhead{Exercice résolu \thepopexr}{poporange}{#1}},
  before lower={\poppill{Solution}{popink}{popcream}\par\vspace{6pt}}}
% Exercice (non résolu en place) — la correction est dans la partie Corrigés
\newtcolorbox{exercice}[1][]{popbase,colframe=popink,
  before upper={\stepcounter{popexo}\popboxhead{Exercice \thepopexo}{popink}{#1}}}
% Corrigé
\newtcolorbox{corrige}[1][]{popbase,colframe=popgreen,colback=popgreenL,
  before upper={\popboxhead{Corrigé}{popgreen}{#1}}}
% Objectifs / prérequis / bilan
\newtcolorbox{objectifs}{popbase,colframe=popink,colback=poporangeL,
  before upper={\popboxhead{Objectifs de la leçon}{popink}{}}}
\newtcolorbox{prerequis}{popbase,colframe=popink,colback=popblueL,
  before upper={\popboxhead{Prérequis}{popblue}{}}}
\newtcolorbox{bilan}{popbase,colframe=popink,colback=white,boxrule=2.8pt,
  before upper={\popboxhead{Fiche bilan}{poporange}{L'essentiel à connaître pour l'examen}}}
% Figure encadrée avec légende
\newtcolorbox{popfigure}[1][]{enhanced,breakable=false,colback=white,colframe=popline,boxrule=1.4pt,arc=8pt,
  left=10pt,right=10pt,top=10pt,bottom=8pt,before skip=10pt,after skip=12pt,halign=center,
  lower separated=false,halign lower=center,
  fontlower=\popsemi\fontsize{9}{11.5}\selectfont\color{popmuted},
  #1}
\newcommand{\legende}[1]{\par\vspace{6pt}{\popsemi\fontsize{9}{11.5}\selectfont\color{popmuted}#1}}

% ============================================================
% Signature OnBuch+
% ============================================================
\newcommand{\poplogoinline}[1]{{\poplogo\fontsize{#1}{#1}\selectfont\color{popink}OnBuch\color{poporange}+}}
\newcommand{\signatureonbuch}{%
  \begin{tcolorbox}[enhanced,colback=popink,colframe=popink,boxrule=2.2pt,arc=10pt,
      drop shadow=poporange,left=14pt,right=14pt,top=10pt,bottom=10pt,before skip=0pt,after skip=0pt]
    \tikz[baseline=-0.7ex]\node[circle,fill=popgreen,draw=popcream,line width=1.3pt,minimum size=22pt,inner sep=0]{\popcheck{white}};%
    \hspace{9pt}\begin{minipage}[c]{0.62\linewidth}
      {\popxbold\fontsize{12}{13}\selectfont\color{popcream}Signé\ \textbullet\ {\poplogo OnBuch\color{poporange}+}}\par\vspace{2pt}
      {\raggedright\popsemi\fontsize{8}{10.5}\selectfont\color{popline}\MakeUppercase{Équipe pédagogique OnBuch+}\\\MakeUppercase{Conforme au programme officiel MINESEC}\par}
    \end{minipage}\hfill
    {\poplogo\fontsize{22}{22}\selectfont\color{popcream}OnBuch\color{poporange}+}
  \end{tcolorbox}%
}

% ============================================================
% Page de garde
% #1 titre · #2 matière · #3 niveau · #4 module · #5 n° leçon · #6 durée ·
% #7 description / objectifs (texte ou itemize)
% ============================================================
\newcommand{\pagegarde}[7]{%
  \thispagestyle{empty}
  \newgeometry{a4paper,margin=1.7cm,top=1.6cm,bottom=1.4cm}%
  \begin{tikzpicture}[remember picture,overlay]
    \fill[poporange!18] ([xshift=-3.2cm,yshift=-3.6cm]current page.north east) circle (4.6cm);
    \fill[poppurple!14] ([xshift=2.4cm,yshift=7.6cm]current page.south west) circle (3.8cm);
  \end{tikzpicture}%
  {\poplogo\fontsize{30}{30}\selectfont\color{popink}OnBuch\color{poporange}+}\hfill
  \raisebox{0.6ex}{\poppill{Cours signé \textbullet\ Premium}{popink}{popcream}}\par
  \vspace{1pt}\popsquiggle{poppurple}\par
  \vspace{8pt}
  \begin{tcolorbox}[enhanced,colback=poporange,colframe=popink,boxrule=2.8pt,arc=13pt,
      drop shadow=popink,left=18pt,right=18pt,top=12pt,bottom=13pt,after skip=10pt]
    \poppill{#2 \textbullet\ #3}{popink}{popcream}\hspace{5pt}\poppill{#4}{white}{popink}\par\vspace{8pt}
    {\popdisplay\fontsize{14}{14}\selectfont\color{popink}LEÇON #5}\par\vspace{4pt}
    {\raggedright\hyphenpenalty=10000\exhyphenpenalty=10000\popdisplay\fontsize{25}{27}\selectfont\color{popcream}\MakeUppercase{#1}\par}
    \vspace{8pt}
    {\popxbold\fontsize{9.5}{9.5}\selectfont\color{popink}\MakeUppercase{Durée conseillée : #6}}
  \end{tcolorbox}
  \begin{tcolorbox}[enhanced,colback=white,colframe=popink,boxrule=2.2pt,arc=10pt,
      drop shadow=popink,left=16pt,right=16pt,top=10pt,bottom=10pt,after skip=8pt]
    \poppill{Dans cette leçon}{poppurple}{white}\par\vspace{5pt}
    \popsemi\fontsize{9.3}{12.6}\selectfont\color{popink}#7
  \end{tcolorbox}
  \noindent\begin{minipage}[t]{0.487\linewidth}
    \begin{tcolorbox}[enhanced,colback=popgreen,colframe=popink,boxrule=2.2pt,arc=10pt,
        drop shadow=popink,left=11pt,right=11pt,top=10pt,bottom=10pt,height=3.1cm,valign=center]
      \tcbox[colback=white,colframe=popink,boxrule=1.3pt,arc=5pt,boxsep=0pt,nobeforeafter,
        left=3pt,right=3pt,top=3pt,bottom=3pt,valign=center]{\qrcode[height=1.9cm]{https://play.google.com/store/apps/details?id=com.luvvix.onbuch&hl=fr}}%
      \hspace{8pt}\begin{minipage}[c]{0.5\linewidth}
        {\popdisplay\fontsize{11}{12.5}\selectfont\color{white}\MakeUppercase{Télécharge OnBuch}}\par\vspace{4pt}
        {\raggedright\popsemi\fontsize{8.4}{11}\selectfont\color{white}Gratuit \textbullet\ Google Play\\Tuteur IA, annales, concours, résultats\par}
      \end{minipage}
    \end{tcolorbox}
  \end{minipage}\hfill
  \begin{minipage}[t]{0.487\linewidth}
    \begin{tcolorbox}[enhanced,colback=poppurple,colframe=popink,boxrule=2.2pt,arc=10pt,
        drop shadow=popink,left=11pt,right=11pt,top=10pt,bottom=10pt,height=3.1cm,valign=center]
      \tikz[baseline=-0.7ex]\node[circle,fill=white,draw=popink,line width=1.3pt,minimum size=1.6cm,inner sep=0]{\popdisplay\fontsize{24}{24}\selectfont\color{poppurple}+};%
      \hspace{8pt}\begin{minipage}[c]{0.6\linewidth}
        {\popdisplay\fontsize{11}{12.5}\selectfont\color{white}\MakeUppercase{Passe à OnBuch+}}\par\vspace{4pt}
        {\raggedright\popsemi\fontsize{8.4}{11}\selectfont\color{white}Tous les cours signés, Tuteur IA illimité, fiches PDF — onglet OnBuch+ de l'app\par}
      \end{minipage}
    \end{tcolorbox}
  \end{minipage}
  \vspace{8pt}
  \popcontenu
  \vfill
  \signatureonbuch
  \restoregeometry
  \newpage
}

% Rangée « Ce cours contient » (page de garde)
\newcommand{\poptile}[3]{% #1 couleur #2 gros texte #3 légende
  \begin{tcolorbox}[enhanced,colback=#1,colframe=popink,boxrule=2pt,arc=9pt,shadow={2.5pt}{-2.5pt}{0pt}{popink},
      left=6pt,right=6pt,top=9pt,bottom=9pt,halign=center,valign=center,height=1.75cm,width=\linewidth,nobeforeafter]
    {\popdisplay\fontsize{12.5}{13}\selectfont\color{white}\MakeUppercase{#2}}\par\vspace{3pt}
    {\centering\popsemi\fontsize{7.8}{9.5}\selectfont\color{white}#3\par}
  \end{tcolorbox}}
\newcommand{\popcontenu}{%
  \par\noindent{\popxboldls\fontsize{7.5}{7.5}\selectfont\color{popmuted}CE COURS SIGNÉ CONTIENT}\par\vspace{7pt}
  \noindent\begin{minipage}[t]{0.235\linewidth}\poptile{popblue}{Cours}{complet, illustré, pas à pas}\end{minipage}\hfill
  \begin{minipage}[t]{0.235\linewidth}\poptile{poporange}{Méthodes}{et exercices résolus}\end{minipage}\hfill
  \begin{minipage}[t]{0.235\linewidth}\poptile{poppink}{Exercices}{type examen + corrigés}\end{minipage}\hfill
  \begin{minipage}[t]{0.235\linewidth}\poptile{popgreen}{Fiche bilan}{l'essentiel à retenir}\end{minipage}\par}

% Sommaire
\newtcolorbox{sommaire}{enhanced,colback=white,colframe=popink,boxrule=2.2pt,arc=10pt,
  drop shadow=popink,left=18pt,right=18pt,top=13pt,bottom=13pt,before skip=6pt,after skip=20pt,
  before upper={\poppill{Sommaire}{poppurple}{white}\par\vspace{9pt}\popsemi\fontsize{10.6}{14.5}\selectfont\color{popink}}}

% Signature de fin de cours — poussée tout en bas de la dernière page
\newcommand{\findecours}{%
  \par\vfill\nopagebreak
  \begin{center}\popsquiggle{poporange}\end{center}\vspace{4pt}
  \signatureonbuch
}
\AtBeginDocument{\def\llbracket{\mathopen{[\mkern-3.5mu[}}\def\rrbracket{\mathclose{]\mkern-3.5mu]}}} % intervalles d'entiers (glyphe ⟦ absent de la police)
% Glyphes absents de Fira Math : redéfinis à partir de glyphes disponibles
\AtBeginDocument{%
  \def\setminus{\mathbin{\backslash}}\let\smallsetminus\setminus
  \def\vdots{\mathord{\vbox{\baselineskip3.2pt\lineskiplimit0pt\kern4pt\hbox{.}\hbox{.}\hbox{.}}}}%
  \def\ddots{\mathinner{\mkern1mu\raise6.5pt\hbox{.}\mkern2mu\raise3.6pt\hbox{.}\mkern2mu\raise0.7pt\hbox{.}\mkern1mu}}%
  \def\triangle{\mathord{\scalebox{0.9}{$\symup{\Delta}$}}}%
  \def\nabla{\mathord{\raisebox{\depth}{\scalebox{1}[-1]{$\symup{\Delta}$}}}}%
  \def\checkmark{\popcheck{popdark}}%
  \def\star{\mathbin{\ast}}\def\bigstar{\mathord{\ast}}%
  \def\diamond{\mathbin{\scalebox{0.6}{$\Diamond$}}}%
  \def\bigcup{\mathop{\mathchoice{\raisebox{-0.3em}{\scalebox{1.7}{$\cup$}}}{\scalebox{1.25}{$\cup$}}{\cup}{\cup}}\displaylimits}%
  \def\bigcap{\mathop{\mathchoice{\raisebox{-0.3em}{\scalebox{1.7}{$\cap$}}}{\scalebox{1.25}{$\cap$}}{\cap}{\cap}}\displaylimits}%
  \def\lesseqqgtr{\mathrel{\lessgtr}}\def\bigoplus{\mathop{\oplus}\displaylimits}%
}
\providecommand{\ssi}{si et seulement si} % abréviation fréquente
\AtBeginDocument{\providecommand{\wideparen}{\overparen}} % arc de cercle

\begin{document}

\pagegarde{Lesson 22 — Grammar: Other forms of comparing, e.g. double, degree, equality, etc}{Anglais}{Tle A}{Chapter 2 : Economic Life and Occupations}{EN22}{2 h}{Cette leçon de grammaire explore les autres formes de comparaison en anglais : l'égalité (as...as), la proportion (the more...the more), la gradation (more and more) et la multiplication (twice as...as). À travers l'observation d'exemples contextualisés, les élèves construisent les règles, évitent les pièges des francophones et s'entraînent à utiliser ces structures dans des phrases et de courts textes.
\vspace{4pt}
\begin{multicols}{2}\small
\begin{itemize}
\item À la fin de cette leçon, je sais construire la comparaison d'égalité avec as + adjectif/adverbe + as
\item À la fin de cette leçon, je sais exprimer l'inégalité avec not as/so...as et less...than
\item À la fin de cette leçon, je sais exprimer la proportion avec la structure the + comparatif...the + comparatif
\item À la fin de cette leçon, je sais exprimer la gradation avec more and more et les comparatifs répétés
\item À la fin de cette leçon, je sais exprimer le double, le triple avec twice/three times as...as
\item À la fin de cette leçon, je sais comparer des quantités avec as much/as many...as
\end{itemize}\end{multicols}}

\begin{sommaire}
\begin{multicols}{2}
\begin{enumerate}
\item Observation et découverte
\item La comparaison d'égalité et d'inégalité
\item La proportion : the more... the more
\item La gradation et la multiplication
\item Comparaison français/anglais et pièges
\item Entraînement guidé et production
\item Activité d'intégration
\item Méthodes et exercices résolus
\item Exercices
\item Corrigés des exercices
\item Fiche bilan
\end{enumerate}
\end{multicols}
\end{sommaire}

\begin{objectifs}
\begin{itemize}
\item À la fin de cette leçon, je sais construire la comparaison d'égalité avec as + adjectif/adverbe + as
\item À la fin de cette leçon, je sais exprimer l'inégalité avec not as/so...as et less...than
\item À la fin de cette leçon, je sais exprimer la proportion avec la structure the + comparatif...the + comparatif
\item À la fin de cette leçon, je sais exprimer la gradation avec more and more et les comparatifs répétés
\item À la fin de cette leçon, je sais exprimer le double, le triple avec twice/three times as...as
\item À la fin de cette leçon, je sais comparer des quantités avec as much/as many...as
\item À la fin de cette leçon, je sais éviter les erreurs fréquentes des francophones (as...than, more and more + adjectif long, etc.)
\item À la fin de cette leçon, je sais utiliser ces structures dans des phrases et un court paragraphe de comparaison
\end{itemize}
\end{objectifs}

\begin{prerequis}
\begin{itemize}
\item Le comparatif de supériorité (adjectifs courts en -er, adjectifs longs avec more)
\item Le superlatif (the + -est / the most)
\item La distinction adjectifs courts / adjectifs longs
\item Les adverbes dérivés des adjectifs (quick → quickly)
\item Les noms dénombrables et indénombrables (much / many)
\end{itemize}
\end{prerequis}

\newpage

\coursec{Observation et découverte}

Au marché Mokolo à Yaoundé, Amina s'arrête devant deux étals de tissus. Elle compare les prix à voix haute : \eng{“This wax print is twice as expensive as the other one, but the quality is just as good.”} Plus tard, en discutant avec son frère qui étudie à Buea, elle remarque : \eng{“The more you practise English, the easier it becomes.”} Ces phrases de la vie quotidienne camerounaise utilisent des formes de comparaison qui vont bien au-delà du simple comparatif en \eng{-er} ou \eng{more} appris en classe de première. Comment exprimer l'\cle{égalité}, la \cle{proportion} ou le \cle{double} en anglais sans faire de faute ? C'est ce que nous allons découvrir dans cette leçon, en partant — comme toujours — de l'observation d'un texte authentique.

\courssub{Lecture guidée : deux commerçantes au marché de Bamenda}

\begin{methode}[Observer avant de mémoriser : la démarche en quatre étapes]
Face à un texte qui contient une structure grammaticale nouvelle, ne cherchez jamais la règle tout de suite. Procédez ainsi :
\begin{enumerate}
\item \textbf{Soulignez} les structures qui se répètent ou qui vous semblent inhabituelles (ici : les groupes de mots autour de \eng{as}, \eng{the more}, \eng{twice}).
\item \textbf{Identifiez} les mots qui encadrent la structure : quel mot la commence ? Quel mot la termine ? Qu'y a-t-il au milieu (adjectif, adverbe, nom) ?
\item \textbf{Comparez} avec le français : comment diriez-vous la même chose dans votre langue ? Où la construction anglaise diffère-t-elle ?
\item \textbf{Formulez une hypothèse} sur la règle, à l'écrit, avant de lire l'explication du professeur. Une règle devinée puis vérifiée se retient dix fois mieux qu'une règle simplement recopiée.
\end{enumerate}
\end{methode}

Lisez maintenant le texte suivant. Les structures de comparaison y sont volontairement nombreuses : à vous de les repérer.

\begin{textebox}[At Bamenda Main Market — an original text]
Every morning at six, Mama Ngong opens her stall at Bamenda Main Market. She sells \textbf{as many vegetables as} Mama Rose, her neighbour, but her stall is \textbf{twice as busy}. Why? Her tomatoes are \textbf{as fresh as} the ones in the big supermarkets of the town, and her prices are \textbf{not as high as} theirs. “The more customers I serve, the more tired I become,” she laughs, “but \textbf{the more I work, the more money I earn}.”

Business is getting \textbf{better and better} every month. Last year, she earned about \textbf{half as much as} she earns today. Now her income is \textbf{twice as high as} it was two years ago, and her daughter's school fees are no longer a problem. Mama Rose is not jealous: her own stall is doing \textbf{better and better} too, although she has \textbf{fewer customers than} her friend. “Mama Ngong works \textbf{twice as hard as} me,” she admits with a smile. “The earlier you arrive at the market, the more you sell. That is the secret.”

The two women agree on one thing: the market is becoming \textbf{more and more crowded}, and rents are getting \textbf{higher and higher}. “The more the town grows, the more expensive life becomes,” sighs Mama Rose. Still, neither of them would exchange her stall for an office job. “An office salary is not half as interesting as a good day at the market,” says Mama Ngong. “Here, the harder you work, the more you earn. Nothing is as satisfying as that.”
\end{textebox}

\textbf{Glossaire.}

\begin{center}
\begin{tabularx}{\linewidth}{|X|X|X|}\hline
\rowcolor{popblueL} Mot anglais & Nature & Traduction française \\ \hline
stall & nom & étal, stand (au marché) \\ \hline
busy & adjectif & animé, très fréquenté \\ \hline
fresh & adjectif & frais, fraîche \\ \hline
to earn & verbe & gagner (de l'argent) \\ \hline
income & nom & revenu \\ \hline
school fees & nom pluriel & frais de scolarité \\ \hline
jealous & adjectif & jaloux, jalouse \\ \hline
crowded & adjectif & bondé, plein de monde \\ \hline
rent & nom & loyer \\ \hline
to sigh & verbe & soupirer \\ \hline
satisfying & adjectif & satisfaisant, gratifiant \\ \hline
\end{tabularx}
\end{center}

\textbf{Questions de compréhension.}

\begin{enumerate}
\item \eng{Whose stall is busier, Mama Ngong's or Mama Rose's?}
\item \eng{Why do customers prefer Mama Ngong's tomatoes?}
\item \eng{How has Mama Ngong's income changed over the last two years?}
\item \eng{What is “the secret” of success at the market, according to Mama Rose?}
\item \eng{Would the two women like to work in an office? Why (not)?}
\end{enumerate}

\begin{corrige}[Questions de compréhension]
\begin{enumerate}
\item \eng{Mama Ngong's stall is busier — it is twice as busy as Mama Rose's.} « L'étal de Mama Ngong est plus animé — deux fois plus que celui de Mama Rose. »
\item \eng{Because they are as fresh as the ones in the big supermarkets, but her prices are not as high.} « Parce qu'elles sont aussi fraîches que celles des supermarchés, mais à des prix moins élevés. »
\item \eng{It has doubled: it is now twice as high as it was two years ago.} « Il a doublé : il est deux fois plus élevé qu'il y a deux ans. »
\item \eng{The earlier you arrive at the market, the more you sell.} « Plus tôt on arrive au marché, plus on vend. »
\item \eng{No, they would not, because they find market work more satisfying and better paid when one works hard.} « Non, car elles trouvent le travail au marché plus gratifiant et mieux payé quand on travaille dur. »
\end{enumerate}
\end{corrige}

\courssub{Repérage des structures de comparaison}

Relevons maintenant toutes les structures de comparaison du texte et classons-les. C'est l'étape centrale de la découverte : vous devez voir par vous-même qu'il n'existe pas une seule comparaison en anglais, mais \textbf{quatre grandes familles}.

\begin{center}
\begin{tabularx}{\linewidth}{|X|X|X|}\hline
\rowcolor{popblueL} Exemple tiré du texte & Catégorie & Traduction française \\ \hline
\eng{as many vegetables as Mama Rose} & égalité & autant de légumes que Mama Rose \\ \hline
\eng{as fresh as the ones in the supermarkets} & égalité & aussi fraîches que celles des supermarchés \\ \hline
\eng{not as high as theirs} & inégalité & pas aussi élevés que les leurs \\ \hline
\eng{the more I work, the more money I earn} & proportion & plus je travaille, plus je gagne d'argent \\ \hline
\eng{the earlier you arrive, the more you sell} & proportion & plus tôt on arrive, plus on vend \\ \hline
\eng{better and better} & gradation & de mieux en mieux \\ \hline
\eng{more and more crowded} & gradation & de plus en plus bondé \\ \hline
\eng{higher and higher} & gradation & de plus en plus élevés \\ \hline
\eng{twice as busy} & multiplication & deux fois plus animé \\ \hline
\eng{twice as high as it was} & multiplication & deux fois plus élevé qu'avant \\ \hline
\eng{half as much as she earns today} & multiplication & la moitié de ce qu'elle gagne aujourd'hui \\ \hline
\end{tabularx}
\end{center}

\begin{exemplebox}[Premiers exemples traduits]
\eng{She sells as many vegetables as her neighbour.} « Elle vend autant de légumes que sa voisine. »

\eng{Her stall is twice as busy as Mama Rose's.} « Son étal est deux fois plus animé que celui de Mama Rose. »

\eng{The more customers she serves, the more tired she becomes.} « Plus elle sert de clients, plus elle se fatigue. »

\eng{Business is getting better and better every month.} « Les affaires vont de mieux en mieux chaque mois. »
\end{exemplebox}

\courssub{Classement : les quatre familles de comparaison}

\begin{definition}[Les quatre formes de comparaison au-delà du comparatif simple]
\begin{itemize}
\item L'\cle{égalité} (\eng{equality}) : \eng{as + adjectif/adverbe + as} — « aussi... que », et sa variante \eng{as many/much... as} — « autant de... que ».
\item L'\cle{inégalité} (\eng{inequality}) : \eng{not as/so... as} — « pas aussi... que », « moins... que ».
\item La \cle{proportion} (\eng{proportion}) : \eng{the + comparatif..., the + comparatif...} — « plus..., plus... ».
\item La \cle{gradation} (\eng{gradation}) : \eng{comparatif + and + comparatif} — « de plus en plus... », « de mieux en mieux ».
\item La \cle{multiplication} (\eng{multiplication}) : \eng{twice / three times / half + as... as} — « deux fois / trois fois / moitié... que ».
\end{itemize}
\end{definition}

\begin{popfigure}
\begin{tikzpicture}[mindmap, grow cyclic, every node/.style={concept, text=white, align=center, font=\small},
root concept/.append style={concept color=popblue, font=\small\bfseries},
level 1/.append style={level distance=3.6cm, sibling angle=90},
level 2/.append style={level distance=2.4cm, sibling angle=45, font=\scriptsize}]
\node[root concept]{Other forms of comparing}
  child[concept color=popgreen]{ node[align=center]{equality\\ \eng{as... as}}
    child[concept color=popgreen!70]{ node[align=center]{\eng{as fresh as}\\ aussi... que} }
    child[concept color=popgreen!70]{ node[align=center]{\eng{as many as}\\ autant que} }
  }
  child[concept color=poporange]{ node[align=center]{proportion\\ \eng{the more...}\\\eng{the more}}
    child[concept color=poporange!70]{ node[align=center]{\eng{the earlier,}\\\eng{the more}} }
  }
  child[concept color=poppurple]{ node[align=center]{gradation\\ \eng{more and more}}
    child[concept color=poppurple!70]{ node[align=center]{\eng{better and}\\\eng{better}} }
    child[concept color=poppurple!70]{ node[align=center]{\eng{higher and}\\\eng{higher}} }
  }
  child[concept color=poppink]{ node[align=center]{multiplication\\ \eng{twice as... as}}
    child[concept color=poppink!70]{ node[align=center]{\eng{half as}\\\eng{much as}} }
    child[concept color=poppink!70]{ node[align=center]{\eng{three times}\\\eng{as... as}} }
  };
\end{tikzpicture}
\legende{Carte mentale : les quatre familles de structures comparatives étudiées dans cette leçon.}
\end{popfigure}

\courssub{Hypothèses des élèves : deviner la règle avant de la lire}

Avant toute explication du professeur, formulez vos propres hypothèses. Voici les questions qui guident votre raisonnement, avec les réponses qu'un bon élève peut déjà deviner à partir des exemples.

\textbf{Hypothèse 1 — l'égalité.} Dans \eng{as fresh as}, le mot \eng{as} apparaît deux fois et encadre l'adjectif. Hypothèse : la structure est \eng{as + adjectif + as}. Attention, l'adjectif reste à la forme simple : on dit \eng{as fresh as}, jamais \eng{as fresher as}.

\textbf{Hypothèse 2 — la négation.} Dans \eng{not as high as theirs}, le \eng{not} se place devant le premier \eng{as}. Hypothèse : pour dire « moins... que », l'anglais nie simplement la structure d'égalité : \eng{not as... as}. Le français « moins cher que » se dit donc \eng{not as expensive as} ou \eng{cheaper than} — deux possibilités, une seule structure à la fois.

\textbf{Hypothèse 3 — la proportion.} Dans \eng{the more I work, the more money I earn}, chaque partie commence par \eng{the} suivi d'un comparatif. Hypothèse : le français « plus..., plus... » se traduit par \eng{the + comparatif..., the + comparatif...}. Ce \eng{the} n'est pas l'article ordinaire : il ne se traduit pas.

\textbf{Hypothèse 4 — la gradation.} Dans \eng{better and better}, le comparatif est répété avec \eng{and}. Hypothèse : « de plus en plus » se dit \eng{comparatif + and + comparatif} : \eng{more and more crowded}, \eng{higher and higher}.

\textbf{Hypothèse 5 — la multiplication.} Dans \eng{twice as busy}, le multiplicateur se place \emph{avant} la structure d'égalité. Hypothèse : \eng{twice / three times / half + as + adjectif + as}. Le mot \eng{double} existe en anglais, mais c'est un adjectif ou un nom (\eng{a double portion}), pas un adverbe de comparaison : on ne dit jamais \eng{double as busy}.

\begin{exoresolu}[Mise en commun des hypothèses]
Énoncé. Sans connaître encore les règles, complétez les phrases suivantes en vous aidant uniquement des exemples du texte : (1) \eng{My tomatoes are ......... fresh ......... yours.} (égalité) ; (2) \eng{......... harder you work, ......... more you earn.} (proportion) ; (3) \eng{Life in Douala is getting ......... and ......... expensive.} (gradation) ; (4) \eng{Her stall is ......... as big as mine.} (multiplication, ×2).
\tcblower
Solution détaillée. (1) La structure d'égalité encadre l'adjectif avec deux \eng{as} : \eng{My tomatoes are \textbf{as} fresh \textbf{as} yours.} (2) Chaque membre commence par \eng{the} + comparatif : \eng{\textbf{The} harder you work, \textbf{the} more you earn.} (3) L'adjectif long \eng{expensive} prend \eng{more} : \eng{getting \textbf{more} and \textbf{more} expensive}. (4) Le multiplicateur précède le premier \eng{as} : \eng{\textbf{twice} as big as mine}. Ces quatre réponses seront confirmées et approfondies dans les sections suivantes de la leçon.
\end{exoresolu}

\courssub{Traduire pour faire émerger les différences}

La traduction est ici un outil de découverte : elle révèle que le français et l'anglais ne découpent pas la réalité de la même façon.

\begin{center}
\begin{tabularx}{\linewidth}{|X|X|X|}\hline
\rowcolor{popblueL} Français & English & Remarque \\ \hline
aussi grand que & as tall as & deux \eng{as}, adjectif simple \\ \hline
autant de clients que & as many customers as & \eng{many} devant un nom comptable \\ \hline
autant d'argent que & as much money as & \eng{much} devant un nom indénombrable \\ \hline
moins cher que / pas aussi cher que & not as expensive as & l'anglais préfère souvent nier l'égalité \\ \hline
plus..., plus... & the more..., the more... & \eng{the} obligatoire, non traduit \\ \hline
plus tôt..., plus... & the earlier..., the more... & idem avec un adverbe \\ \hline
de mieux en mieux & better and better & comparatif répété avec \eng{and} \\ \hline
de plus en plus cher & more and more expensive & \eng{more} répété pour les adjectifs longs \\ \hline
deux fois plus grand que & twice as big as & le multiplicateur avant \eng{as} \\ \hline
deux fois moins / moitié moins & half as much as & \eng{half} suit le même schéma \\ \hline
\end{tabularx}
\end{center}

\begin{attention}[Pièges déjà visibles pour les francophones]
\begin{itemize}
\item Ne traduisez jamais « plus..., plus... » par \eng{more..., more...} seul : il faut \eng{\textbf{the} more..., \textbf{the} more...}. Ce \eng{the} est obligatoire même si le français n'a pas d'article.
\item « Deux fois plus » ne se dit pas \eng{two times more} dans la structure d'égalité : dites \eng{\textbf{twice} as big as}. (\eng{Two times} se comprend, mais \eng{twice} est la forme correcte et naturelle.)
\item Dans \eng{as... as}, l'adjectif ne prend jamais la marque du comparatif : \eng{as \textbf{good} as}, pas \eng{as better as}.
\item « De plus en plus » avec un adjectif court : \eng{bigger and bigger}, pas \eng{more and more big}. Avec un adjectif long : \eng{more and more difficult}.
\item Faux ami : \eng{actually} ne signifie pas « actuellement » mais « en fait » ; dans les comparaisons, méfiez-vous aussi de \eng{to compare with} (comparer à, pour établir un parallèle) et \eng{compared to} (par rapport à).
\end{itemize}
\end{attention}

\begin{experience}[À vous de jouer : observation en binômes]
Par groupes de deux, relisez le texte du marché de Bamenda à voix haute. L'élève A lit, l'élève B lève la main à chaque structure de comparaison et annonce sa catégorie (\eng{equality}, \eng{proportion}, \eng{gradation} ou \eng{multiplication}). Inversez ensuite les rôles. Terminez en rédigeant ensemble deux phrases sur votre propre quartier ou votre marché, une avec \eng{as... as} et une avec \eng{the more..., the more...}. Exemple attendu : \eng{The more customers a trader has, the more money she makes.}
\end{experience}

\begin{exercice}[Premier entraînement : classez et traduisez]
\begin{enumerate}
\item Classez les structures suivantes dans la bonne famille (égalité, inégalité, proportion, gradation, multiplication) : \eng{three times as expensive} ; \eng{not so difficult as} ; \eng{worse and worse} ; \eng{the sooner, the better} ; \eng{as much time as}.
\item Traduisez en anglais : « Yaoundé est deux fois plus grand que Bafoussam. » ; « Plus tu révises, meilleurs sont tes résultats. » ; « Le transport devient de plus en plus cher. »
\end{enumerate}
\end{exercice}

\begin{corrige}[Premier entraînement]
\begin{enumerate}
\item \eng{three times as expensive} : multiplication ; \eng{not so difficult as} : inégalité ; \eng{worse and worse} : gradation ; \eng{the sooner, the better} : proportion ; \eng{as much time as} : égalité.
\item \eng{Yaoundé is twice as big as Bafoussam.} ; \eng{The more you revise, the better your results are.} ; \eng{Transport is getting more and more expensive.}
\end{enumerate}
\end{corrige}

\begin{aretenir}[Ce qu'il faut retenir de cette observation]
L'anglais dispose de quatre structures de comparaison au-delà du comparatif simple : l'égalité \eng{as... as}, la proportion \eng{the more..., the more...}, la gradation \eng{better and better} et la multiplication \eng{twice as... as}. Chacune a une construction fixe qu'il faut observer, deviner, puis vérifier. Le français traduit ces structures par « aussi... que », « plus..., plus... », « de plus en plus » et « deux fois plus... que » — mais les mécanismes internes (place du \eng{the}, place du multiplicateur, forme de l'adjectif) sont différents. C'est précisément ce décalage que les sections suivantes vont systématiser.
\end{aretenir}

\coursec{La comparaison d'égalité et d'inégalité}

Dans la section précédente, nous avons observé un texte et un dialogue dans lesquels apparaissaient des structures comme \eng{as cheap as}, \eng{not as crowded as} ou encore \eng{as much money as}. Nous avons remarqué que ces formes ne se contentent pas de dire qu'une chose est « plus » ou « moins » qu'une autre : elles expriment aussi l'\cle{égalité} (\eng{as...as}) et l'\cle{inégalité} par le bas (\eng{not as...as}, \eng{less...than}). Il est temps maintenant de dégager les règles précises de construction et d'emploi de ces structures, qui sont au cœur de l'expression de la comparaison en anglais — et une source d'erreurs très fréquente chez les francophones, tant à l'épreuve de grammaire qu'à l'épreuve de composition du Baccalauréat.

\courssub{Observer avant de conclure : un texte support}

Lisons d'abord ce court article original, rédigé dans le style de la presse économique anglophone. Soulignez mentalement toutes les structures de comparaison que vous repérez : il y en a plus de dix.

\begin{textebox}[Two Markets, One City — adapted from a business report]
Mokolo Market and Mvog-Mbi Market are both famous in Yaoundé, but they are not equally busy. Mokolo is \textbf{just as large as} Mvog-Mbi, yet it attracts far more customers every morning. A trader can earn \textbf{as much money as} a shopkeeper in the city centre, especially during the festive season. However, the stalls at Mvog-Mbi are \textbf{not as crowded as} those at Mokolo, and the car park is \textbf{nearly as big as} the one at the central market. Prices for tomatoes are \textbf{almost as low as} in the villages around the city, but imported goods are \textbf{less affordable than} local products. Interestingly, Mvog-Mbi has \textbf{as many taxi stops as} Mokolo, although the roads leading to it are \textbf{not as wide as} the main avenue. A young seller told our reporter: “My tomatoes sell \textbf{as fast as} my neighbour's, but I work \textbf{less comfortably than} she does because my stall is smaller.” For many families, shopping at Mvog-Mbi is \textbf{as convenient as} going to a supermarket, and often cheaper. Economists note that small markets are \textbf{as important as} large companies for the local economy: they employ \textbf{as many young people as} some factories do. In short, the two markets are equal in size, but unequal in energy — and that difference is exactly what makes each of them unique.
\end{textebox}

\textbf{Glossaire.}

\begin{center}
\begin{tabularx}{\linewidth}{|l|l|X|}
\hline
\rowcolor{popblueL} Mot anglais & Nature & Traduction française \\
\hline
stall & nom & étal, stand (de marché) \\
\hline
affordable & adjectif & abordable (financièrement) \\
\hline
convenient & adjectif & pratique, commode \\
\hline
to employ & verbe & employer, donner du travail à \\
\hline
unequal & adjectif & inégal \\
\hline
\end{tabularx}
\end{center}

\textbf{Questions de compréhension.}
\begin{enumerate}
\item Which market is more crowded in the morning?
\item Are the two markets the same size? Justify with a sentence from the text.
\item What do small markets and large companies have in common, according to economists?
\end{enumerate}

\textbf{Éléments de réponse.} 1. \eng{Mokolo is more crowded: its stalls are busier than those at Mvog-Mbi.} 2. \eng{Yes, they are: “Mokolo is just as large as Mvog-Mbi.”} 3. \eng{They are equally important for the local economy: small markets employ as many young people as some factories do.}

\courssub{La comparaison d'égalité : as + adjectif/adverbe + as}

Reprenons les exemples du texte : \eng{just as large as}, \eng{as convenient as}, \eng{as fast as}, \eng{as important as}. Dans tous les cas, l'adjectif ou l'adverbe reste à la \cle{forme simple} (jamais de \eng{-er}, jamais de \eng{more}), et il est encadré par deux mots : \eng{as} avant, \eng{as} après.

\begin{propriete}[Règle : la comparaison d'égalité]
Pour dire que deux éléments sont \textbf{égaux} sur un point, on utilise la structure :
\begin{center}
\textbf{as + adjectif / adverbe (forme simple) + as + terme de comparaison}
\end{center}
\begin{itemize}
\item Avec un \textbf{adjectif} : \eng{Douala is as hot as Garoua.} « Douala est aussi chaude que Garoua. »
\item Avec un \textbf{adverbe} : \eng{She speaks English as fluently as her teacher.} « Elle parle anglais aussi couramment que son professeur. »
\end{itemize}
L'adjectif ou l'adverbe ne prend \textbf{jamais} de marque de comparaison : on dit \eng{as tall as}, jamais \emph{as taller as}.
\end{propriete}

\begin{exemplebox}[Exemples traduits — l'égalité]
\begin{itemize}
\item \eng{Douala is as hot as Garoua.} « Douala est aussi chaude que Garoua. »
\item \eng{My sister is as tall as my mother.} « Ma sœur est aussi grande que ma mère. »
\item \eng{This exercise is as difficult as the last one.} « Cet exercice est aussi difficile que le précédent. »
\item \eng{The bendskin driver rides as fast as a taxi.} « Le moto-taximan roule aussi vite qu'un taxi. »
\item \eng{Our team played as well as the champions.} « Notre équipe a joué aussi bien que les champions. »
\end{itemize}
\end{exemplebox}

\begin{attention}[Piège n° 1 : jamais « as...than »]
Le français dit « aussi... \textbf{que} ». Beaucoup d'élèves traduisent ce « que » par \eng{than} et écrivent \emph{as tall than}, ce qui est \textbf{faux}. Retenez :
\begin{itemize}
\item égalité : \eng{as...as} — \eng{He is as tall as his father.} « Il est aussi grand que son père. »
\item supériorité : \eng{more/-er...than} — \eng{He is taller than his father.} « Il est plus grand que son père. »
\end{itemize}
\eng{than} ne se combine \textbf{jamais} avec \eng{as}.
\end{attention}

\begin{aretenir}[Une structure voisine : the same...as]
Pour exprimer l'identité, l'anglais utilise aussi \eng{the same + nom + as} :
\begin{itemize}
\item \eng{She goes to the same school as her cousin.} « Elle va à la même école que sa cousine. »
\item \eng{My phone is the same as yours.} « Mon téléphone est le même que le vôtre. »
\end{itemize}
Attention à la préposition : c'est \eng{the same \textbf{as}}, jamais \emph{the same than} ni \emph{the same like}. À l'opposé : \eng{different \textbf{from}} — \eng{My job is different from yours.} « Mon métier est différent du vôtre. »
\end{aretenir}

\courssub{La comparaison d'inégalité : not as/so...as et less...than}

Pour exprimer l'\cle{infériorité} (« moins... que »), l'anglais dispose de deux structures.

\begin{propriete}[Règle : la comparaison d'inégalité]
\textbf{1. La négation de l'égalité} (la plus fréquente) :
\begin{center}
\textbf{not as + adjectif/adverbe + as} \quad ou \quad \textbf{not so + adjectif/adverbe + as}
\end{center}
\eng{This taxi is not as fast as the bus.} « Ce taxi n'est pas aussi rapide que le bus. »\\
\eng{Bamenda is not so crowded as Douala.} « Bamenda n'est pas aussi peuplée que Douala. »

\textbf{Remarque importante :} \eng{so} ne peut remplacer le premier \eng{as} \textbf{que dans une phrase négative}. On ne dit jamais \emph{so tall as} dans une phrase affirmative.

\textbf{2. La structure avec less} (infériorité directe) :
\begin{center}
\textbf{less + adjectif/adverbe + than}
\end{center}
\eng{Imported goods are less affordable than local products.} « Les produits importés sont moins abordables que les produits locaux. »\\
\eng{I work less comfortably than she does.} « Je travaille moins confortablement qu'elle. »
\end{propriete}

\begin{exemplebox}[Exemples traduits — l'inégalité]
\begin{itemize}
\item \eng{This taxi is not as fast as the bus.} « Ce taxi n'est pas aussi rapide que le bus. »
\item \eng{The second exercise is not so difficult as the first.} « Le deuxième exercice n'est pas aussi difficile que le premier. »
\item \eng{The village road is less busy than the main avenue.} « La route du village est moins fréquentée que l'avenue principale. »
\item \eng{He speaks less clearly than his brother.} « Il parle moins clairement que son frère. »
\end{itemize}
\end{exemplebox}

\begin{attention}[Piège n° 2 : less et les adjectifs courts]
Avec les adjectifs \textbf{courts} (une syllabe), l'anglais préfère souvent \eng{not as...as} à \eng{less...than}. \eng{less big than} est grammaticalement possible mais peu naturel ; dites plutôt \eng{not as big as}. En revanche, avec les adjectifs \textbf{longs} (\eng{expensive, comfortable, important}), \eng{less...than} est tout à fait correct : \eng{This phone is less expensive than mine.} « Ce téléphone est moins cher que le mien. »
\end{attention}

\courssub{La comparaison des quantités : as much...as / as many...as}

Le texte support contenait deux phrases essentielles : \eng{a trader can earn as much money as a shopkeeper} et \eng{Mvog-Mbi has as many taxi stops as Mokolo}. Pourquoi \eng{much} dans un cas et \eng{many} dans l'autre ? Parce que \eng{money} est \cle{indénombrable} et \eng{taxi stops} est \cle{dénombrable}.

\begin{propriete}[Règle : comparer les quantités]
\begin{center}
\textbf{as much + nom indénombrable + as} \qquad \textbf{as many + nom dénombrable pluriel + as}
\end{center}
\begin{itemize}
\item \eng{He earns as much money as his brother.} « Il gagne autant d'argent que son frère. »
\item \eng{She drinks as much water as her coach recommends.} « Elle boit autant d'eau que son entraîneur le recommande. »
\item \eng{Our school has as many students as the lycée downtown.} « Notre école a autant d'élèves que le lycée du centre-ville. »
\item \eng{I have read as many books as you.} « J'ai lu autant de livres que toi. »
\end{itemize}
À la forme négative : \eng{I don't have as much time as you.} « Je n'ai pas autant de temps que toi. »
\end{propriete}

\begin{attention}[Piège n° 3 : much ou many ?]
Comme pour \eng{much/many} en général, le choix dépend de la nature du nom :
\begin{itemize}
\item \textbf{indénombrables} (\eng{money, water, time, rice, information, work}) $\rightarrow$ \eng{as much...as} ;
\item \textbf{dénombrables au pluriel} (\eng{books, students, taxis, chairs}) $\rightarrow$ \eng{as many...as}.
\end{itemize}
Ne dites jamais \emph{as many money as} ni \emph{as much books as}. Rappel : \eng{information} et \eng{news} sont indénombrables en anglais, contrairement au français « informations » et « nouvelles ».
\end{attention}

\courssub{Nuance et renforcement : just as, nearly as, almost as}

L'anglais module finement l'égalité grâce à des adverbes placés \textbf{avant} le premier \eng{as} :

\begin{center}
\begin{tabularx}{\linewidth}{|l|X|X|}
\hline
\rowcolor{popblueL} Structure & Exemple anglais & Traduction \\
\hline
\eng{just as...as} & \eng{Mokolo is just as large as Mvog-Mbi.} & Mokolo est exactement aussi grand que Mvog-Mbi. \\
\hline
\eng{nearly as...as} & \eng{The car park is nearly as big as the other one.} & Le parking est presque aussi grand que l'autre. \\
\hline
\eng{almost as...as} & \eng{Prices are almost as low as in the villages.} & Les prix sont presque aussi bas qu'aux villages. \\
\hline
\eng{not nearly as...as} & \eng{The bus is not nearly as fast as the train.} & Le bus est loin d'être aussi rapide que le train. \\
\hline
\eng{twice as...as} & \eng{Douala is twice as big as Buea.} & Douala est deux fois plus grande que Buea. \\
\hline
\end{tabularx}
\end{center}

\courssub{Tableau récapitulatif : égalité et inégalité}

\begin{popfigure}
\begin{tikzpicture}[
  box/.style={draw=popline, rounded corners=3pt, fill=#1, text width=6.4cm, align=center, inner sep=6pt, font=\small},
  head/.style={draw=popblue, fill=popblue, text=white, rounded corners=3pt, text width=6.4cm, align=center, font=\small\bfseries, inner sep=5pt}
]
\node[head] (h1) at (0,0) {Égalité};
\node[head, right=0.6cm of h1] (h2) {Inégalité (infériorité)};
\node[box=popgreenL, below=0.25cm of h1, align=center] (a1){\textbf{as + adj. + as}\\ \eng{as tall as} \\ « aussi grand que »};
\node[box=poppinkL, below=0.25cm of h2, align=center] (b1){\textbf{not as/so + adj. + as}\\ \eng{not as expensive as} \\ « pas aussi cher que »};
\node[box=popgreenL, below=0.2cm of a1, align=center] (a2){\textbf{as much + indénombrable + as}\\ \eng{as much money as} \\ « autant d'argent que »};
\node[box=poppinkL, below=0.2cm of b1, align=center] (b2){\textbf{less + adj. + than}\\ \eng{less crowded than} \\ « moins bondé que »};
\node[box=popgreenL, below=0.2cm of a2, align=center] (a3){\textbf{as many + dénombrable + as}\\ \eng{as many books as} \\ « autant de livres que »};
\node[box=poppinkL, below=0.2cm of b2, align=center] (b3){\textbf{not as many/much + nom + as}\\ \eng{not as many taxis as} \\ « pas autant de taxis que »};
\end{tikzpicture}
\legende{Les structures de l'égalité et de l'inégalité en un coup d'œil}
\end{popfigure}

\begin{center}
\begin{tabularx}{\linewidth}{|X|X|X|}
\hline
\rowcolor{popblueL} Français & English & Remarque \\
\hline
aussi grand que & as tall as & adjectif à la forme simple \\
\hline
aussi vite que & as fast as & \eng{fast} est un adverbe irrégulier (pas de \eng{-ly}) \\
\hline
aussi bien que & as well as & adverbe irrégulier de \eng{good} \\
\hline
pas aussi loin que & not as far as & \eng{so} possible : \eng{not so far as} \\
\hline
autant d'argent que & as much money as & indénombrable \\
\hline
autant de clients que & as many customers as & dénombrable \\
\hline
moins cher que & less expensive than & adjectif long \\
\hline
exactement aussi... que & just as...as & renforcement de l'égalité \\
\hline
le même... que & the same...as & préposition \eng{as}, jamais \eng{than} \\
\hline
\end{tabularx}
\end{center}

\begin{attention}[Piège n° 4 : les adverbes irréguliers]
Certains adverbes ne se forment pas avec \eng{-ly}. Dans la comparaison d'égalité, on utilise leur forme propre :
\begin{itemize}
\item \eng{good} (adj.) $\rightarrow$ \eng{well} (adv.) : \eng{She sings as well as her sister.} « Elle chante aussi bien que sa sœur. »
\item \eng{fast} reste \eng{fast} : \eng{He runs as fast as a bendskin.} « Il court aussi vite qu'un bendskin. » (jamais \emph{fastly})
\item \eng{hard} reste \eng{hard} : \eng{They work as hard as their parents.} « Ils travaillent aussi dur que leurs parents. » (\eng{hardly} signifie « à peine » !)
\end{itemize}
\end{attention}

\begin{savaistu}[Le saviez-vous ?]
La structure \eng{as...as} est si productive en anglais qu'elle a donné naissance à des centaines d'expressions figées, souvent imagées : \eng{as busy as a bee} « occupé comme une abeille », \eng{as cold as ice} « froid comme la glace », \eng{as old as the hills} « vieux comme le monde », \eng{as good as gold} « sage comme une image ». Les reconnaître vous aidera à mieux comprendre les textes littéraires et journalistiques du Baccalauréat — et à enrichir vos compositions.
\end{savaistu}

\courssub{Entraînement : exercice résolu}

\begin{exoresolu}[Compléter avec la bonne structure]
Complétez chaque phrase avec \eng{as...as}, \eng{not as...as}, \eng{as much...as}, \eng{as many...as} ou \eng{less...than}.

1. A kilo of cocoa costs \ldots\ money \ldots\ a kilo of imported rice. (égalité)\\
2. The evening bus is \ldots\ crowded \ldots\ the morning bus. (infériorité, négation)\\
3. There are \ldots\ passengers \ldots\ seats on this bus. (égalité, dénombrable)\\
4. My new phone is \ldots\ expensive \ldots\ my old one. (infériorité avec \eng{less})\\
5. She types \ldots\ quickly \ldots\ a professional secretary. (égalité, adverbe)

\tcblower
\textbf{Solution détaillée.}

1. \eng{A kilo of cocoa costs as much money as a kilo of imported rice.} — \eng{money} est indénombrable, donc \eng{as much...as}. « Un kilo de cacao coûte autant d'argent qu'un kilo de riz importé. »

2. \eng{The evening bus is not as crowded as the morning bus.} — négation de l'égalité : \eng{not as + adjectif + as}. « Le bus du soir n'est pas aussi bondé que celui du matin. »

3. \eng{There are as many passengers as seats on this bus.} — \eng{passengers} est dénombrable, donc \eng{as many...as}. « Il y a autant de passagers que de sièges dans ce bus. »

4. \eng{My new phone is less expensive than my old one.} — infériorité avec \eng{less + adjectif long + than}. « Mon nouveau téléphone est moins cher que mon ancien. »

5. \eng{She types as quickly as a professional secretary.} — égalité avec un adverbe régulier en \eng{-ly}. « Elle tape aussi vite qu'une secrétaire professionnelle. »
\end{exoresolu}

\begin{exercice}[À vous de jouer]
Traduisez en anglais :
\begin{enumerate}
\item Yaoundé est presque aussi grande que Douala.
\item Ce marché n'est pas aussi animé que le marché central.
\item Il gagne autant d'argent que son frère.
\item Nous avons moins de devoirs que les élèves de première.
\item Elle parle anglais aussi bien que sa cousine.
\end{enumerate}
\end{exercice}

\begin{corrige}[Exercice « À vous de jouer »]
\begin{enumerate}
\item \eng{Yaoundé is almost as big as Douala.}
\item \eng{This market is not as lively as the central market.}
\item \eng{He earns as much money as his brother.}
\item \eng{We have less homework than the Form Five students.} (ou : \eng{We do not have as much homework as the Form Five students.})
\item \eng{She speaks English as well as her cousin.}
\end{enumerate}
\end{corrige}

Retenons l'essentiel : l'égalité se construit avec \eng{as + adjectif/adverbe + as}, l'inégalité avec \eng{not as/so...as} ou \eng{less...than}, et les quantités avec \eng{as much...as} (indénombrable) ou \eng{as many...as} (dénombrable). Dans la prochaine section, nous verrons comment l'anglais exprime la \cle{proportion} avec la célèbre structure \eng{the more...the more}.

\coursec{La proportion : the more... the more}

Dans la section précédente, nous avons comparé deux éléments entre eux : \eng{John is as tall as Peter} (égalité), \eng{Yaoundé is not as hot as Douala} (inégalité). Mais il arrive qu'on ne compare pas deux personnes ou deux choses fixes : on veut exprimer que \textbf{deux réalités varient ensemble}, l'une entraînant l'autre. En français, on dit : « \emph{Plus} tu travailles, \emph{plus} tu gagnes d'argent. » En anglais, il existe une structure parallèle très élégante : \eng{The more you work, the more money you earn.} C'est la \cle{comparaison de proportion}, objet de cette section.

\courssub{Observation : découvrir la structure}

Lisons d'abord ce court texte. Soulignez mentalement toutes les structures répétées.

\begin{textebox}[A trader's wisdom — texte d'observation]
Mama Etoundi sells second-hand clothes at the Mokolo market in Yaoundé. She often gives advice to her young niece who has just opened a small shop. “Listen, my daughter,” she says. “The earlier you arrive at the market, the better places you find. The harder you work, the more customers you attract. And the more customers you have, the more money you make — it is as simple as that!”

Her niece smiles and answers: “And the more money I make, the happier my family will be!”

Mama Etoundi laughs. “Exactly. But remember: the higher your prices, the fewer customers you will keep. Business is a balance. The more you listen to people, the more they trust you. The more they trust you, the longer they stay with you.”

That evening, the niece writes in her notebook: “The more I learn, the more confident I become. The more confident I become, the more successful my shop will be.” She closes the notebook and whispers: “The sooner I start, the better.”
\end{textebox}

\textbf{Glossaire :} \emph{second-hand} (adj., d'occasion) ; \emph{a place} (n., un emplacement) ; \emph{to attract} (v., attirer) ; \emph{to trust} (v., faire confiance à) ; \emph{a balance} (n., un équilibre) ; \emph{confident} (adj., sûr de soi) ; \emph{to whisper} (v., chuchoter).

\textbf{Questions d'observation :}
\begin{enumerate}
\item Combien de fois la structure \eng{the + comparatif} apparaît-elle dans le texte ?
\item Chaque phrase contient-elle une ou deux propositions ?
\item Où se trouve le comparatif par rapport au sujet : avant ou après ?
\end{enumerate}

Vous avez remarqué la régularité parfaite : \textbf{chaque phrase contient deux propositions, et chacune commence par \eng{the} + un comparatif}. C'est toute la mécanique de la structure. Remarquez aussi, dans le texte, la correction idiomatique : on dit \eng{the \textbf{higher} your prices} (« plus tes prix sont élevés ») et non \eng{the more expensive your prices} — en anglais, ce sont les \emph{choses} qui sont \emph{expensive}, les \emph{prix} sont \emph{high} ou \emph{low}.

\courssub{La règle : forme et sens}

\begin{propriete}[La comparaison de proportion]
Pour exprimer que deux choses \textbf{varient ensemble proportionnellement} (quand l'une augmente, l'autre augmente ou diminue), l'anglais utilise la structure :

\begin{center}
\textbf{The + comparatif + sujet + verbe, \quad the + comparatif + sujet + verbe.}
\end{center}

\begin{itemize}
\item \eng{The more you study, the better your results are.} — Plus tu étudies, meilleurs sont tes résultats.
\item \eng{The less you spend, the more you save.} — Moins tu dépenses, plus tu économises.
\end{itemize}

\textbf{Points essentiels :}
\begin{enumerate}
\item Les \textbf{deux} propositions commencent par \eng{the} + comparatif. Ce \eng{the} est obligatoire : ce n'est pas l'article ordinaire, mais un mot figé de la construction.
\item Le comparatif \textbf{précède le sujet} dans les deux propositions : \eng{the harder \textbf{you} work} et non \eng{you work harder}.
\item La virgule sépare les deux propositions.
\end{enumerate}
\end{propriete}

\begin{attention}[Ne pas confondre avec l'article !]
Dans \eng{the more you work}, le mot \eng{the} ne veut rien dire tout seul : il ne désigne rien de précis. Il fait partie d'un bloc insécable \eng{the + comparatif}. C'est l'équivalent exact du « plus... plus » français. Ne cherchez donc pas à le traduire par « le/la ».
\end{attention}

\courssub{Adjectifs courts et adjectifs longs}

La structure fonctionne avec tous les comparatifs, qu'ils soient courts (en \emph{-er}) ou longs (avec \emph{more}).

\begin{center}
\begin{tabularx}{\linewidth}{|X|X|X|}
\hline
\rowcolor{popblueL} Type de comparatif & Exemple anglais & Traduction \\
\hline
Court (\emph{-er}) : hard $\rightarrow$ harder & The harder you work, the more you earn. & Plus tu travailles dur, plus tu gagnes. \\
\hline
Court : early $\rightarrow$ earlier & The earlier you book, the cheaper the ticket. & Plus tôt tu réserves, moins le billet est cher. \\
\hline
Court : soon $\rightarrow$ sooner & The sooner we leave, the sooner we arrive. & Plus tôt nous partons, plus tôt nous arrivons. \\
\hline
Long (\emph{more} + adj.) & The more expensive the product, the more profitable the business. & Plus le produit est cher, plus l'affaire est rentable. \\
\hline
Long : more difficult & The more difficult the exam, the harder the students concentrate. & Plus l'examen est difficile, plus les élèves se concentrent. \\
\hline
Irrégulier : good $\rightarrow$ better & The more you practise, the better you speak. & Plus tu pratiques, mieux tu parles. \\
\hline
Irrégulier : bad $\rightarrow$ worse & The less you sleep, the worse you feel. & Moins tu dors, plus mal tu te sens. \\
\hline
\emph{less} (indénombrables) & The less you eat, the thinner you get. & Moins tu manges, plus tu maigris. \\
\hline
\emph{fewer} (dénombrables) & The fewer mistakes you make, the higher your mark. & Moins tu fais de fautes, plus ta note est élevée. \\
\hline
\end{tabularx}
\end{center}

\begin{exemplebox}[Exemples traduits — à apprendre par cœur]
\begin{itemize}
\item \eng{The more you practise, the more fluent you become.} — Plus tu pratiques, plus tu deviens à l'aise.
\item \eng{The earlier you book, the cheaper the ticket.} — Plus tôt tu réserves, moins le billet est cher.
\item \eng{The more cocoa Cameroon exports, the more money the country earns.} — Plus le Cameroun exporte de cacao, plus le pays gagne d'argent.
\item \eng{The older I get, the more patient I become.} — Plus je vieillis, plus je deviens patient.
\item \eng{The faster you drive, the more dangerous the journey.} — Plus tu conduis vite, plus le voyage est dangereux.
\item \eng{The more we discuss, the better we understand each other.} — Plus nous discutons, mieux nous nous comprenons.
\end{itemize}
\end{exemplebox}

\courssub{L'ordre des mots : le point délicat}

C'est ici que les francophones trébuchent le plus souvent. En français, on dit : « plus \textbf{tu} \textbf{étudies} » — le sujet vient tout de suite après « plus ». En anglais, le comparatif passe \textbf{avant} le sujet, et le reste de la proposition suit l'ordre normal sujet + verbe.

\begin{center}
\begin{tabularx}{\linewidth}{|X|X|X|}
\hline
\rowcolor{popblueL} Français & English & Remarque \\
\hline
Plus tu étudies, plus tu réussis. & The more you study, the more you succeed. & Comparatif avant le sujet dans les deux propositions. \\
\hline
Plus le magasin est grand, plus il y a de clients. & The bigger the shop, the more customers it has. & Pas d'inversion sujet-verbe : \eng{it has}, jamais \eng{has it}. \\
\hline
Moins tu parles, plus tu écoutes. & The less you talk, the more you listen. & \eng{less} suit exactement la même règle. \\
\hline
\end{tabularx}
\end{center}

\begin{attention}[Pas d'inversion interrogative !]
Même si la phrase commence par un groupe en tête, il n'y a \textbf{jamais d'inversion} du sujet et du verbe : on dit \eng{The more you earn, the more tax \textbf{you pay}} et non \eng{the more tax do you pay}. L'ordre reste celui d'une phrase affirmative ordinaire.
\end{attention}

\begin{popfigure}
\begin{tikzpicture}[font=\small]
% Flèche 1 : cause
\draw[-{Stealth[length=4mm]}, line width=2.2pt, popblue] (0,0) -- (4.6,2.1);
\node[align=center, text=popblue, above=2pt] at (2.3,1.05) {The more you study};
% Flèche 2 : effet
\draw[-{Stealth[length=4mm]}, line width=2.2pt, poporange] (6.4,0) -- (11,2.1);
\node[align=center, text=poporange, above=2pt] at (8.7,1.05) {the better your results};
% Lien cause-effet
\draw[-{Stealth[length=3mm]}, line width=1.4pt, popgreen, dashed] (4.9,1.05) -- (6.1,1.05);
\node[align=center, text=popgreen, below=4pt] at (5.5,1.05) {proportion};
% Étiquettes
\node[align=center, text=popmuted] at (2.3,-0.5) {cause : action qui augmente};
\node[align=center, text=popmuted] at (8.7,-0.5) {effet : résultat qui suit};
\end{tikzpicture}
\legende{La proportion : quand la première quantité monte, la seconde monte avec elle.}
\end{popfigure}

\courssub{La forme raccourcie : proverbes et conversation}

Quand le sens est évident, l'anglais \textbf{supprime le sujet et le verbe} et ne garde que les deux comparatifs. C'est la forme préférée des proverbes et de la conversation quotidienne.

\begin{exemplebox}[Formes raccourcies célèbres]
\begin{itemize}
\item \eng{The more, the merrier.} — Plus on est de fous, plus on rit. (littéralement : plus il y en a, plus c'est joyeux)
\item \eng{The sooner, the better.} — Le plus tôt sera le mieux.
\item \eng{The bigger, the better.} — Plus c'est grand, mieux c'est.
\item \eng{The less said, the better.} — Moins on en parle, mieux c'est.
\item \eng{The cheaper, the better.} — Moins c'est cher, mieux c'est.
\end{itemize}
\end{exemplebox}

\begin{savaistu}[“The more, the merrier”]
L'expression \eng{The more, the merrier} est extrêmement courante dans la conversation britannique. On l'emploie pour inviter chaleureusement quelqu'un à se joindre à un groupe : \eng{“Can I bring my cousin to the party?” — “Of course! The more, the merrier!”} (« Je peux amener mon cousin à la fête ? — Bien sûr ! Plus on est de fous, plus on rit ! »). Le mot \eng{merry} (« joyeux ») est le même que dans \eng{Merry Christmas}. Retenez ce proverbe : il vous fera gagner des points à l'oral comme à l'écrit.
\end{savaistu}

\courssub{La proportion à tous les temps}

La structure se conjugue sans difficulté au présent, au passé et au futur. Seuls les verbes changent ; l'ossature \eng{the + comparatif} reste identique.

\begin{center}
\begin{tabular}{|l|l|}
\hline
\rowcolor{popblueL} Temps & Exemple \\
\hline
Présent & The more he works, the more he earns. \\
\hline
Passé & The more he worked, the more he earned. \\
\hline
Présent + futur (\emph{will}) & The more he works, the more he will earn. \\
\hline
Futur dans les deux propositions & The more you will practise, the better you will play. \\
\hline
Présent + futur (conseil) & The earlier you start, the sooner you will finish. \\
\hline
\end{tabular}
\end{center}

\begin{propriete}[Concordance des temps]
\begin{itemize}
\item \textbf{Présent + présent} : vérité générale, habitude. \eng{The more sugar you eat, the more weight you gain.}
\item \textbf{Passé + passé} : récit au passé. \eng{The faster we ran, the more tired we became.}
\item \textbf{Présent + futur} : prédiction (la deuxième proposition prend souvent \eng{will}). \eng{The more you save now, the easier your life will be.}
\end{itemize}
\end{propriete}

\begin{exemplebox}[Exemples au passé et au futur — contexte camerounais]
\begin{itemize}
\item \eng{The more the farmers planted, the more cocoa they harvested last season.} — Plus les planteurs ont planté, plus ils ont récolté de cacao la saison dernière.
\item \eng{The more you invest in your education, the more opportunities you will have.} — Plus tu investis dans tes études, plus tu auras d'opportunités.
\item \eng{The harder the team trained, the better they played in Buea.} — Plus l'équipe s'est entraînée dur, mieux elle a joué à Buea.
\item \eng{The more bendskins there are in Douala, the more careful pedestrians will have to be.} — Plus il y a de bendskins à Douala, plus les piétons devront être prudents.
\end{itemize}
\end{exemplebox}

\courssub{Exercice résolu et pièges}

\begin{exoresolu}[Traduction et construction]
\textbf{Énoncé.} Traduisez en anglais :
\begin{enumerate}
\item Plus tu lis, plus ton vocabulaire s'enrichit.
\item Moins nous dépensons, plus nous économiserons.
\item Plus le match était difficile, plus les supporters criaient.
\item Plus c'est simple, mieux c'est.
\end{enumerate}
\tcblower
\textbf{Solution détaillée.}
\begin{enumerate}
\item \eng{The more you read, the more your vocabulary grows.} — Comparatif \eng{more} placé avant le sujet \eng{you} ; deuxième proposition au présent (vérité générale).
\item \eng{The less we spend, the more we will save.} — \eng{less} pour l'infériorité ; futur \eng{will save} dans la seconde proposition car il s'agit d'une prédiction.
\item \eng{The more difficult the match was, the louder the supporters shouted.} — Adjectif long : \eng{the more difficult} ; passé dans les deux propositions (\eng{was}, \eng{shouted}) ; \eng{louder} = « plus fort » (comparatif court de \eng{loud}).
\item \eng{The simpler, the better.} — Forme raccourcie : sujet et verbe supprimés, comme dans un proverbe.
\end{enumerate}
\end{exoresolu}

\begin{attention}[Les quatre erreurs classiques des francophones]
\begin{enumerate}
\item \textbf{Oublier le \eng{the}.} Ne dites jamais \eng{More you work, more you earn.} Le français « plus... plus » se traduit par \eng{\textbf{the} more... \textbf{the} more}. Sans \eng{the}, la phrase est fautive.
\item \textbf{Traduire « plus... plus » par \eng{more... more} seuls.} Même erreur : \eng{more} tout seul signifie « davantage », pas « plus... plus ».
\item \textbf{Inverser sujet et verbe.} Pas de \eng{the more do you work} : l'ordre est \eng{the more \textbf{you work}}.
\item \textbf{Utiliser \eng{less} avec un dénombrable.} On dit \eng{the \textbf{fewer} customers} (clients, dénombrable) mais \eng{the \textbf{less} money} (argent, indénombrable).
\end{enumerate}
\end{attention}

\begin{exercice}[À vous de jouer]
\begin{enumerate}
\item Complétez : The \dots (long) you wait, the \dots (difficult) it becomes.
\item Mettez dans l'ordre : earn / the / you / more / work / the / harder / you / more / the.
\item Traduisez : « Plus on s'entraîne, plus on progresse. »
\item Traduisez : « Moins tu dors, plus tu seras fatigué. »
\item Trouvez la forme raccourcie : « Le plus tôt sera le mieux. »
\end{enumerate}
\end{exercice}

\begin{corrige}[Exercice « À vous de jouer »]
\begin{enumerate}
\item \eng{The \textbf{longer} you wait, the \textbf{more difficult} it becomes.} — \eng{long} est un adjectif court (\emph{-er}) ; \eng{difficult} est long (\emph{more}).
\item \eng{The harder you work, the more you earn.}
\item \eng{The more you train, the more you progress.} (ou \eng{the better you perform}).
\item \eng{The less you sleep, the more tired you will be.} — Futur dans la seconde proposition (prédiction).
\item \eng{The sooner, the better.}
\end{enumerate}
\end{corrige}

\begin{aretenir}[L'essentiel de la section]
\begin{itemize}
\item La proportion s'exprime par : \eng{\textbf{The} + comparatif + sujet + verbe, \textbf{the} + comparatif + sujet + verbe.}
\item Le comparatif précède toujours le sujet ; jamais d'inversion.
\item Adjectifs courts (\eng{the harder}), longs (\eng{the more expensive}), irréguliers (\eng{the better}, \eng{the worse}) et \eng{the less} / \eng{the fewer} : tous suivent la même structure.
\item Forme raccourcie très fréquente : \eng{The more, the merrier.} / \eng{The sooner, the better.}
\item La structure fonctionne au présent, au passé et au futur.
\item Piège n° 1 : « plus... plus » = \eng{the more... the more}, jamais \eng{more... more} seuls.
\end{itemize}
\end{aretenir}

\coursec{La gradation et la multiplication}

Dans la section précédente, nous avons vu comment exprimer une relation de proportion entre deux évolutions avec la structure \eng{the more... the more...} : \eng{The more you practise, the better you speak.} Cette structure met en relation \textbf{deux} phénomènes qui progressent ensemble. Mais comment décrire l'évolution \textbf{d'un seul} phénomène, qui s'intensifie peu à peu dans le temps ? Et comment comparer des quantités chiffrées avec précision (\emph{deux fois plus}, \emph{trois fois plus}, \emph{moitié moins}) ? C'est l'objet de cette section : la \cle{gradation} (\eng{better and better}, \eng{more and more difficult}) et la \cle{multiplication} (\eng{twice as big as}, \eng{three times as many students as}).

\courssub{Observation : un texte pour découvrir}

\begin{textebox}[Economic life in Douala — reading text]
Life in Douala is changing fast. The city is growing \textbf{bigger and bigger}: new buildings appear every month, and the roads are becoming \textbf{more and more crowded}. For many families, daily life is getting \textbf{harder and harder} because prices keep rising. A bag of rice that cost 15,000 francs CFA two years ago now costs 30,000 francs: it is \textbf{twice as expensive as} before. Transport is not better. A bendskin ride across the city can cost \textbf{three times as much as} it did five years ago.

Yet there is good news too. The port of Douala handles \textbf{twice as many containers as} it did a decade ago, and more and more young people are starting small businesses. “The situation is becoming \textbf{more and more difficult},” says Mama Eyenga, a tomato seller at Marché Central, “but we are working \textbf{harder and harder}, and our children are getting \textbf{better and better} results at school. That gives me hope.”

Economists say the city needs \textbf{twice as much investment as} it receives today. Without it, traffic jams will become \textbf{longer and longer}, and the cost of living will rise \textbf{higher and higher}.
\end{textebox}

\begin{definition}[Glossaire du texte]
\begin{itemize}
\item \eng{crowded} (adj.) : bondé, encombré
\item \eng{a bendskin} (n.) : une moto-taxi (Cameroun)
\item \eng{a decade} (n.) : une décennie
\item \eng{a seller} (n.) : un vendeur, une vendeuse
\item \eng{investment} (n.) : un investissement
\item \eng{a traffic jam} (n.) : un embouteillage
\item \eng{to handle} (v.) : traiter, gérer
\end{itemize}
\end{definition}

\textbf{Questions d'observation.} Observez les mots en gras dans le texte :
\begin{enumerate}
\item Quels adjectifs sont simplement \emph{répétés} avec \eng{and} ? (\eng{bigger and bigger}, \eng{harder and harder}...)
\item Quels adjectifs sont précédés de \eng{more and more} ? (\eng{more and more crowded}, \eng{more and more difficult}...)
\item Quelle différence constatez-vous entre ces deux groupes d'adjectifs ?
\item Dans \eng{twice as expensive as} et \eng{three times as much as}, quel mot revient toujours deux fois ?
\end{enumerate}

\courssub{La gradation : décrire une évolution progressive}

\begin{propriete}[Règle de la gradation — forme et emploi]
Pour exprimer qu'une qualité \textbf{s'intensifie progressivement dans le temps} (« de plus en plus... », « de mieux en mieux... »), l'anglais utilise deux structures, selon la \textbf{longueur de l'adjectif} :
\begin{itemize}
\item \textbf{Adjectif court} (une syllabe, ou deux syllabes en \emph{-y}) : \cle{comparatif + and + comparatif}.\\ \eng{cold} $\rightarrow$ \eng{colder and colder} ; \eng{bad} $\rightarrow$ \eng{worse and worse}.
\item \textbf{Adjectif long} (deux syllabes et plus, hors \emph{-y}) : \cle{more and more + adjectif}.\\ \eng{difficult} $\rightarrow$ \eng{more and more difficult} ; \eng{important} $\rightarrow$ \eng{more and more important}.
\end{itemize}
La gradation est souvent introduite par les verbes \eng{to get} ou \eng{to become} (= devenir) : \eng{It is getting darker and darker.}
\end{propriete}

\textbf{Remarque essentielle.} Le choix entre les deux structures suit \emph{exactement la même règle} que la formation du comparatif ordinaire : les adjectifs qui prennent \emph{-er} au comparatif se répètent avec \eng{and} ; ceux qui prennent \eng{more} utilisent \eng{more and more}. Si vous savez former le comparatif, vous savez former la gradation.

\begin{exemplebox}[Exemples traduits]
\begin{itemize}
\item \eng{The traffic in Yaoundé is getting worse and worse.} = La circulation à Yaoundé devient de pire en pire.
\item \eng{English is becoming more and more important.} = L'anglais devient de plus en plus important.
\item \eng{The nights are getting colder and colder in Bamenda.} = Les nuits deviennent de plus en plus froides à Bamenda.
\item \eng{Your English is getting better and better.} = Ton anglais devient de mieux en mieux.
\item \eng{Smartphones are becoming more and more expensive.} = Les smartphones deviennent de plus en plus chers.
\item \eng{The days are getting shorter and shorter.} = Les jours deviennent de plus en plus courts.
\end{itemize}
\end{exemplebox}

\begin{center}
\begin{tabularx}{\linewidth}{|X|X|X|}\hline
\rowcolor{popblueL} Adjectif & Gradation & Traduction \\ \hline
\eng{good} (court, irrégulier) & \eng{better and better} & de mieux en mieux \\ \hline
\eng{bad} (court, irrégulier) & \eng{worse and worse} & de pire en pire \\ \hline
\eng{big} (court) & \eng{bigger and bigger} & de plus en plus grand \\ \hline
\eng{easy} (deux syllabes en -y) & \eng{easier and easier} & de plus en plus facile \\ \hline
\eng{difficult} (long) & \eng{more and more difficult} & de plus en plus difficile \\ \hline
\eng{crowded} (long) & \eng{more and more crowded} & de plus en plus bondé \\ \hline
\end{tabularx}
\end{center}

\begin{popfigure}
\begin{tikzpicture}[font=\small]
% Frise ascendante : gradation
\node[draw=popblue, fill=popblueL, rounded corners, minimum width=2.6cm, minimum height=0.9cm, align=center] (a) at (0,0) {\eng{good}};
\node[draw=poporange, fill=poporangeL, rounded corners, minimum width=3.2cm, minimum height=0.9cm, align=center] (b) at (4.2,1.1) {\eng{better}\\ \eng{and better}};
\node[draw=popgreen, fill=popgreenL, rounded corners, minimum width=2.6cm, minimum height=0.9cm, align=center] (c) at (8.4,2.2) {\eng{the best}};
\draw[-{Stealth[length=3mm]}, very thick, popdark] (a) -- (b);
\draw[-{Stealth[length=3mm]}, very thick, popdark] (b) -- (c);
\node[popmuted, align=center] at (4.2,-0.9) {gradation : évolution progressive};
% Barres de multiplication
\fill[popblue!60] (10.6,-0.9) rectangle (11.6,0.1);
\node[align=center, anchor=south] at (11.1,0.15) {\eng{1x price}};
\fill[poporange!70] (12.4,-0.9) rectangle (13.4,1.1);
\node[align=center, anchor=south] at (12.9,1.15) {\eng{2x = twice}\\ \eng{as expensive}};
\draw[popdark] (10.4,-0.9) -- (13.6,-0.9);
\end{tikzpicture}
\legende{À gauche : la gradation ascendante (\eng{good} $\rightarrow$ \eng{better and better} $\rightarrow$ \eng{the best}). À droite : la multiplication — une barre deux fois plus haute = \eng{twice as expensive}.}
\end{popfigure}

\courssub{La multiplication : comparer des quantités chiffrées}

\begin{propriete}[Règle de la multiplication — forme et emploi]
Pour exprimer qu'une quantité est \textbf{deux fois, trois fois... supérieure (ou inférieure)} à une autre, l'anglais utilise :
\begin{itemize}
\item \cle{twice / three times / half + as + adjectif + as} :\\ \eng{This bag is twice as heavy as that one.} (Ce sac est deux fois plus lourd que celui-là.)
\item Avec les \textbf{noms} : \cle{twice / three times + as much + nom indénombrable + as} ou \cle{twice / three times + as many + nom dénombrable pluriel + as} :\\ \eng{He earns twice as much money as his brother.}\\ \eng{Our school has three times as many students as yours.}
\end{itemize}
Points de vigilance : \eng{twice} (= deux fois) ne prend jamais de \emph{-s} ; on dit \eng{three times}, \eng{four times}... avec un \emph{-s} ; \eng{half as... as} exprime « moitié moins / deux fois moins » : \eng{The journey takes half as long by taxi.} (Le trajet dure deux fois moins longtemps en taxi.)
\end{propriete}

\begin{exemplebox}[Exemples traduits]
\begin{itemize}
\item \eng{Petrol is twice as expensive as last year.} = L'essence coûte deux fois plus cher que l'année dernière.
\item \eng{Douala is three times as big as Buea.} = Douala est trois fois plus grande que Buea.
\item \eng{She earns twice as much money as her colleague.} = Elle gagne deux fois plus d'argent que son collègue.
\item \eng{There are twice as many girls as boys in this class.} = Il y a deux fois plus de filles que de garçons dans cette classe.
\item \eng{This exercise is half as difficult as the last one.} = Cet exercice est deux fois moins difficile que le précédent.
\end{itemize}
\end{exemplebox}

\begin{center}
\begin{tabularx}{\linewidth}{|X|X|X|}\hline
\rowcolor{popgreenL} Français & English & Remarque \\ \hline
deux fois plus grand que & \eng{twice as big as} & jamais \emph{two times} pour « deux fois » \\ \hline
trois fois plus cher que & \eng{three times as expensive as} & \emph{times} avec un -s dès trois \\ \hline
deux fois plus d'argent que & \eng{twice as much money as} & \eng{much} + indénombrable \\ \hline
trois fois plus d'élèves que & \eng{three times as many students as} & \eng{many} + dénombrable \\ \hline
deux fois moins long que & \eng{half as long as} & \eng{half} = moitié \\ \hline
\end{tabularx}
\end{center}

\courssub{Exercice résolu et pièges à éviter}

\begin{exoresolu}[Application guidée]
\textbf{Énoncé.} Complétez avec la gradation ou la multiplication correcte de l'adjectif entre parenthèses.
\begin{enumerate}
\item The harmattan wind is getting \dots\ (dry).
\item Finding a job is becoming \dots\ (difficult) for young graduates.
\item A taxi ride costs \dots\ (expensive) a bendskin ride. (x3)
\item My grandfather's farm produces \dots\ cocoa \dots\ mine. (much, x2)
\end{enumerate}
\tcblower
\textbf{Solution détaillée.}
\begin{enumerate}
\item \eng{drier and drier} — adjectif court en \emph{-y} (\eng{dry} $\rightarrow$ \eng{drier}) : on répète le comparatif avec \eng{and}. « Le vent de l'harmattan devient de plus en plus sec. »
\item \eng{more and more difficult} — adjectif long (trois syllabes) : \eng{more and more} + adjectif. « Trouver un emploi devient de plus en plus difficile pour les jeunes diplômés. »
\item \eng{three times as expensive as} — multiplication par trois : \eng{three times} + \eng{as} + adjectif + \eng{as}. « Une course en taxi coûte trois fois plus cher qu'une course en moto-taxi. »
\item \eng{twice as much cocoa as} — \eng{cocoa} est indénombrable, donc \eng{much} (et non \eng{many}) ; « deux fois » = \eng{twice}. « La ferme de mon grand-père produit deux fois plus de cacao que la mienne. »
\end{enumerate}
\end{exoresolu}

\begin{attention}[Erreurs fréquentes des francophones]
\begin{itemize}
\item Ne dites \textbf{jamais} \eng{more and more better} ! On choisit \emph{soit} la forme en \emph{-er} répétée (\eng{better and better}), \emph{soit} \eng{more and more} + adjectif long (\eng{more and more interesting}) — jamais les deux ensemble.
\item « Deux fois » se dit \eng{twice}, pas \emph{two times} : \eng{twice as big as} (et non \emph{two times as big as}).
\item N'oubliez pas le \textbf{deuxième} \eng{as} : le français dit « deux fois plus cher \emph{que} », l'anglais dit \eng{twice as expensive \emph{as}} — la structure est symétrique : \eng{as ... as}.
\item \eng{much} ou \eng{many} ? Comme toujours : \eng{much} devant un indénombrable (\eng{money}, \eng{water}, \eng{time}), \eng{many} devant un dénombrable pluriel (\eng{students}, \eng{cars}).
\item Attention à l'orthographe des comparatifs répétés : \eng{big} $\rightarrow$ \eng{bi\textbf{gg}er and bigger} (doublement de la consonne), \eng{easy} $\rightarrow$ \eng{easier and easier} (\emph{y} $\rightarrow$ \emph{i}).
\end{itemize}
\end{attention}

\begin{exercice}[À vous de jouer]
\begin{enumerate}
\item Traduisez : « La vie à Douala devient de plus en plus chère. »
\item Traduisez : « Ce marché est deux fois plus grand que le nôtre. »
\item Complétez : \eng{The rainy season is getting \dots\ (long) and the roads are getting \dots\ (bad).}
\item Complétez : \eng{Yaoundé receives \dots\ (many, x3) visitors \dots\ my village.}
\item Transformez avec \eng{half as ... as} : \eng{A bus ticket costs 500 francs. A taxi ride costs 1,000 francs.}
\end{enumerate}
\end{exercice}

\begin{corrige}[Exercice « À vous de jouer »]
\begin{enumerate}
\item \eng{Life in Douala is becoming more and more expensive.} (adjectif long $\rightarrow$ \eng{more and more}.)
\item \eng{This market is twice as big as ours.} (\eng{twice}, pas \emph{two times} ; \eng{ours} = le nôtre.)
\item \eng{The rainy season is getting longer and longer and the roads are getting worse and worse.} (\eng{bad} a un comparatif irrégulier : \eng{worse}.)
\item \eng{Yaoundé receives three times as many visitors as my village.} (\eng{visitors} est dénombrable $\rightarrow$ \eng{many}.)
\item \eng{A bus ticket is half as expensive as a taxi ride.} (Un ticket de bus coûte deux fois moins cher qu'une course en taxi.)
\end{enumerate}
\end{corrige}

\begin{aretenir}[L'essentiel de la section]
\begin{itemize}
\item \textbf{Gradation} (évolution progressive) : adjectif court $\rightarrow$ \eng{comparatif + and + comparatif} (\eng{worse and worse}) ; adjectif long $\rightarrow$ \eng{more and more + adjectif} (\eng{more and more important}).
\item \textbf{Multiplication} (quantités chiffrées) : \eng{twice / three times / half + as + adjectif + as} ; avec un nom : \eng{twice as much + indénombrable}, \eng{three times as many + dénombrable}.
\item Jamais de mélange : pas de \eng{more and more better} ; « deux fois » = \eng{twice}.
\end{itemize}
\end{aretenir}

\coursec{Comparaison français/anglais et pièges}

Nous venons de voir comment exprimer la gradation (\eng{more and more expensive}) et la multiplication (\eng{twice as much as}). Ces structures sont simples en anglais, mais elles deviennent de véritables pièges dès que l'on traduit mot à mot depuis le français. C'est exactement ce que le correcteur du Bac surveille : l'élève francophone qui pense en français et écrit en anglais. Cette section dresse le bilan contrastif des deux langues et vous donne les réflexes de vérification indispensables.

\courssub{Le tableau contrastif : quatre équivalences à mémoriser}

Commençons par observer quatre phrases authentiques, tirées d'un dialogue entre deux commerçants du marché de Mokolo à Yaoundé :

\begin{exemplebox}[Observation : quatre structures, quatre équivalences]
\begin{itemize}
\item \eng{Plantains are as expensive as cocoyams this season.} = Les bananes plantains sont \textbf{aussi} chères \textbf{que} les macabos cette saison.
\item \eng{The more customers you have, the more money you make.} = \textbf{Plus} tu as de clients, \textbf{plus} tu gagnes d'argent.
\item \eng{Imported rice is getting more and more expensive.} = Le riz importé devient \textbf{de plus en plus} cher.
\item \eng{A taxi ride costs twice as much as a bendskin ride.} = Une course en taxi coûte \textbf{deux fois plus} cher qu'une course en bendskin.
\end{itemize}
\end{exemplebox}

L'observation est claire : là où le français utilise des mots simples (\emph{aussi}, \emph{plus}, \emph{de plus en plus}, \emph{deux fois plus}), l'anglais impose des \cle{structures fixes}, avec des mots-outils obligatoires (\eng{as... as}, \eng{the... the}, \eng{twice}) qu'il est interdit de modifier ou d'omettre.

\begin{center}
\begin{tabularx}{\linewidth}{|X|X|X|}
\hline
\rowcolor{popblueL} \textbf{Français} & \textbf{English} & \textbf{Remarque} \\
\hline
aussi... que & as... as & jamais \eng{as... than} \\
\hline
plus... plus & the more... the more & les deux \eng{the} sont obligatoires \\
\hline
de plus en plus & more and more + adjectif long ; adjectif court + -er and -er & \eng{more and more expensive} ; \eng{bigger and bigger} \\
\hline
deux fois plus... (que) & twice as... as & forme préférée au Bac ; \eng{two times bigger than} est à éviter \\
\hline
trois fois plus... & three times as... as & \eng{three times as much as} \\
\hline
autant de... que & as much... as (indénombrable) ; as many... as (dénombrable) & choisir selon la nature du nom \\
\hline
\end{tabularx}
\end{center}

\begin{popfigure}
\begin{center}
\begin{tikzpicture}
\node[draw=popblue, fill=popblueL, rounded corners, font=\bfseries, align=center] (fr) at (0,0) {Français};
\node[draw=popgreen, fill=popgreenL, rounded corners, font=\bfseries, align=center] (en) at (9,0) {English};
\node[align=center, font=\small] (f1) at (0,-1.4) {aussi grand que};
\node[align=center, font=\small] (e1) at (9,-1.4) {as big as};
\node[align=center, font=\small] (f2) at (0,-2.6) {plus... plus};
\node[align=center, font=\small] (e2) at (9,-2.6) {the more... the more};
\node[align=center, font=\small] (f3) at (0,-3.8) {de plus en plus cher};
\node[align=center, font=\small] (e3) at (9,-3.8) {more and more expensive};
\node[align=center, font=\small] (f4) at (0,-5) {deux fois plus grand};
\node[align=center, font=\small] (e4) at (9,-5) {twice as big as};
\draw[-{Stealth}, thick, popgreen] (f1) -- (e1);
\draw[-{Stealth}, thick, popgreen] (f2) -- (e2);
\draw[-{Stealth}, thick, popgreen] (f3) -- (e3);
\draw[-{Stealth}, thick, popgreen] (f4) -- (e4);
\node[align=center, font=\small\itshape, text=red] (x1) at (4.5,-1.9) {as big than ✗};
\node[align=center, font=\small\itshape, text=red] (x2) at (4.5,-3.1) {more you work, more you earn ✗};
\node[align=center, font=\small\itshape, text=red] (x3) at (4.5,-4.3) {more and more better ✗};
\node[align=center, font=\small\itshape, text=red] (x4) at (4.5,-5.5) {two times bigger than ✗};
\end{tikzpicture}
\end{center}
\legende{Les quatre équivalences clés français/anglais et les formes fautives barrées en rouge.}
\end{popfigure}

\courssub{Piège n°1 : \eng{as... than} au lieu de \eng{as... as}}

C'est l'erreur la plus fréquente chez les candidats francophones. En français, la comparaison de supériorité utilise \emph{que} (\emph{plus grand que}) ; par contamination, l'élève écrit \eng{as big than}. Or \eng{than} n'apparaît en anglais qu'après un \cle{comparatif} (\eng{bigger than}, \eng{more expensive than}). Dans la comparaison d'égalité, la structure est symétrique : \eng{as} + adjectif + \eng{as}.

\begin{attention}[Erreur fréquente des francophones]
\eng{My house is as big than yours.} ✗

\eng{My house is as big as yours.} ✓ = « Ma maison est aussi grande que la tienne. »

Réflexe : si l'adjectif n'est \textbf{pas} au comparatif, on ne peut pas utiliser \eng{than}. On dit \eng{bigger than yours} (comparatif + \eng{than}) ou \eng{as big as yours} (égalité + \eng{as... as}), mais jamais un mélange des deux.
\end{attention}

\courssub{Piège n°2 : oublier le \eng{the} dans \eng{the more... the more}}

En français, \emph{plus tu travailles, plus tu gagnes} ne contient aucun article. L'élève francophone traduit donc spontanément sans \eng{the}. Erreur : la structure anglaise de proportion exige \eng{the} devant \textbf{chaque} comparatif, car il s'agit historiquement d'un ancien instrumental (\emph{by how much more... by so much more}).

\begin{attention}[Erreur fréquente des francophones]
\eng{More you work, more you earn.} ✗

\eng{The more you work, the more you earn.} ✓ = « Plus tu travailles, plus tu gagnes. »

Autres exemples : \eng{The earlier you book, the cheaper the ticket is.} (Plus tu réserves tôt, moins le billet est cher.) — \eng{The more you practise speaking, the more confident you become.} (Plus tu t'entraînes à parler, plus tu deviens confiant.)
\end{attention}

\courssub{Piège n°3 : \eng{more and more better}, la double marque de comparaison}

Les adjectifs courts prennent le suffixe \eng{-er} ; les adjectifs longs prennent \eng{more}. Il est interdit de cumuler les deux marques. Cette règle, déjà vue pour le comparatif simple, s'applique aussi à la gradation : on ne dit jamais \eng{more and more better}, car \eng{better} est déjà le comparatif irrégulier de \eng{good}.

\begin{propriete}[Gradation : une seule marque de comparaison]
\begin{itemize}
\item Adjectif court : on répète le comparatif en \eng{-er} : \eng{The nights are getting shorter and shorter.} = « Les nuits raccourcissent de plus en plus. »
\item Adjectif long : \eng{more and more} + adjectif : \eng{Life in Douala is becoming more and more expensive.} = « La vie à Douala devient de plus en plus chère. »
\item Comparatifs irréguliers : \eng{better and better} (de mieux en mieux), \eng{worse and worse} (de pire en pire).
\end{itemize}
\end{propriete}

\begin{attention}[Double marque interdite]
\eng{The results are getting more and more better.} ✗

\eng{The results are getting better and better.} ✓ = « Les résultats s'améliorent de mieux en mieux. »

De même : \eng{more and more easier} ✗ → \eng{easier and easier} ✓.
\end{attention}

\courssub{Piège n°4 : \eng{as much as} ou \eng{as many as} ?}

Le français dit \emph{autant de} dans tous les cas (\emph{autant d'argent}, \emph{autant d'élèves}). L'anglais, lui, distingue selon la nature du nom :

\begin{propriete}[Le choix much/many dans les comparaisons de quantité]
\begin{itemize}
\item \eng{as much as} + nom \textbf{indénombrable} : \eng{She earns as much money as her brother.} = « Elle gagne autant d'argent que son frère. »
\item \eng{as many as} + nom \textbf{dénombrable pluriel} : \eng{Our school has as many students as yours.} = « Notre école a autant d'élèves que la vôtre. »
\end{itemize}
La même distinction s'applique avec la multiplication : \eng{twice as much water}, \eng{twice as many books}.
\end{propriete}

\begin{attention}[Erreur fréquente des francophones]
\eng{There are as much candidates as last year.} ✗

\eng{There are as many candidates as last year.} ✓ = « Il y a autant de candidats que l'année dernière. »

Test rapide : peut-on mettre un nombre devant le nom ? \emph{Three candidates} oui → dénombrable → \eng{many}. \emph{Three moneys} non → indénombrable → \eng{much}.
\end{attention}

\courssub{Piège n°5 : \eng{two times bigger than} au lieu de \eng{twice as big as}}

Pour traduire \emph{deux fois plus grand}, la tentation est grande de calquer le français : \eng{two times bigger than}. Cette forme, bien qu'entendue parfois à l'oral familier, est considérée comme incorrecte ou maladroite à l'écrit, et elle est sanctionnée au Bac. La forme attendue est \eng{twice as} + adjectif + \eng{as}. Pour trois fois et au-delà, on utilise \eng{three times as... as}, \eng{four times as... as}, etc.

\begin{exemplebox}[Exemples traduits : la multiplication]
\begin{itemize}
\item \eng{This bag costs three times as much as that one.} = « Ce sac coûte trois fois plus cher que celui-là. »
\item \eng{The new market is twice as big as the old one.} = « Le nouveau marché est deux fois plus grand que l'ancien. »
\item \eng{Yaoundé receives twice as much rain as Maroua.} = « Yaoundé reçoit deux fois plus de pluie que Maroua. »
\item \eng{A litre of petrol now costs almost twice as much as five years ago.} = « Un litre d'essence coûte maintenant presque deux fois plus cher qu'il y a cinq ans. »
\end{itemize}
\end{exemplebox}

\begin{attention}[Forme préférée au Bac]
\eng{His salary is two times higher than mine.} ✗ (à éviter à l'écrit)

\eng{His salary is twice as high as mine.} ✓ = « Son salaire est deux fois plus élevé que le mien. »

Retenez aussi que \eng{twice} ne se dit jamais \eng{two times} dans cette structure, et que l'adjectif reste à la base (\eng{as big as}), pas au comparatif (\eng{as bigger as} ✗).
\end{attention}

\courssub{Faux ami : \eng{actually} ne signifie pas \emph{actuellement}}

Dans les comparaisons, on veut souvent corriger une première impression : \emph{En fait, le marché est plus grand que je ne pensais.} Le mot anglais pour \emph{en fait / en réalité} est \eng{actually} — et c'est bien lui qu'il faut utiliser ici, mais attention : \eng{actually} ne signifie \textbf{jamais} \emph{actuellement}. Pour \emph{actuellement}, on dit \eng{currently}, \eng{at present} ou \eng{at the moment}.

\begin{center}
\begin{tabularx}{\linewidth}{|X|X|X|}
\hline
\rowcolor{popblueL} \textbf{Français} & \textbf{English} & \textbf{Exemple} \\
\hline
en fait, en réalité & actually & \eng{The exam was actually easier than I expected.} = « En fait, l'examen était plus facile que je ne le pensais. » \\
\hline
actuellement, en ce moment & currently / at present & \eng{Prices are currently higher than last year.} = « Les prix sont actuellement plus élevés que l'année dernière. » \\
\hline
\end{tabularx}
\end{center}

\begin{attention}[Faux ami]
\eng{Prices are actually higher than last year.} signifie « Les prix sont en fait plus élevés que l'année dernière » (correction d'une idée reçue), et non « Les prix sont actuellement plus élevés ». Pour exprimer le moment présent, écrivez : \eng{Prices are currently higher than last year.}
\end{attention}

\courssub{Méthode de vérification et entraînement}

\begin{methode}[Vérifier sa phrase comparative en trois étapes]
\begin{enumerate}
\item \textbf{Identifier le type de comparaison} : s'agit-il d'une égalité (\emph{aussi... que}), d'une proportion (\emph{plus... plus}), d'une gradation (\emph{de plus en plus}) ou d'une multiplication (\emph{deux fois plus}) ?
\item \textbf{Appliquer la structure exacte} : \eng{as + adj + as} ; \eng{the + comparatif, the + comparatif} ; \eng{more and more + adj} ou \eng{adj-er and adj-er} ; \eng{twice/three times as + adj + as}. Vérifier qu'aucune marque n'est doublée et qu'aucun \eng{the} ne manque.
\item \textbf{Vérifier much/many selon le nom} : si la comparaison porte sur une quantité, demander si le nom est dénombrable (\eng{as many as}) ou indénombrable (\eng{as much as}).
\end{enumerate}
\end{methode}

\begin{exoresolu}[Correction de phrases fautives]
\textbf{Énoncé.} Les phrases suivantes contiennent chacune une erreur typique de francophone. Corrigez-les et justifiez.

1. \eng{My village is as far than yours from the main road.}

2. \eng{More you save, more you can invest.}

3. \eng{The situation is getting more and more worse.}

4. \eng{We don't have as much customers as the supermarket.}

5. \eng{Her shop is two times bigger than mine.}

\tcblower
\textbf{Solution détaillée.}

1. \eng{My village is as far \textbf{as} yours from the main road.} — Comparaison d'égalité : la structure est \eng{as... as} ; \eng{than} n'apparaît qu'après un comparatif (\eng{farther than}).

2. \eng{\textbf{The} more you save, \textbf{the} more you can invest.} — Proportion : chaque comparatif doit être précédé de \eng{the}.

3. \eng{The situation is getting \textbf{worse and worse}.} — \eng{worse} est déjà le comparatif irrégulier de \eng{bad} ; ajouter \eng{more} crée une double marque interdite.

4. \eng{We don't have as \textbf{many} customers as the supermarket.} — \eng{customers} est dénombrable pluriel ; \eng{much} est réservé aux indénombrables.

5. \eng{Her shop is \textbf{twice as big as} mine.} — Pour \emph{deux fois plus}, la forme correcte à l'écrit est \eng{twice as + adjectif + as}, avec l'adjectif à la base et non au comparatif.
\end{exoresolu}

\begin{exercice}[À vous de jouer]
Traduisez en anglais, en appliquant la méthode en trois étapes :

1. Le cacao est aussi important que le café pour notre économie.

2. Plus on exporte, plus le pays gagne d'argent.

3. Le transport devient de plus en plus cher à Bamenda.

4. Un sac de riz coûte trois fois plus cher qu'un sac de maïs.

5. Nous produisons autant de tomates que le village voisin.
\end{exercice}

\begin{corrige}[Exercice : traductions]
1. \eng{Cocoa is as important as coffee for our economy.} (égalité : \eng{as... as})

2. \eng{The more we export, the more money the country earns.} (proportion : \eng{the more... the more})

3. \eng{Transport is getting more and more expensive in Bamenda.} (gradation, adjectif long : \eng{more and more} + adjectif)

4. \eng{A bag of rice costs three times as much as a bag of maize.} (multiplication : \eng{three times as much as}, car le prix est une quantité indénombrable)

5. \eng{We produce as many tomatoes as the neighbouring village.} (\eng{tomatoes} est dénombrable : \eng{as many as})
\end{corrige}

\begin{aretenir}[L'essentiel de la section]
\begin{itemize}
\item Quatre équivalences fixes : \emph{aussi... que} = \eng{as... as} ; \emph{plus... plus} = \eng{the more... the more} ; \emph{de plus en plus} = \eng{more and more} / \eng{-er and -er} ; \emph{deux fois plus} = \eng{twice as... as}.
\item Jamais \eng{as... than}, jamais de proportion sans \eng{the}, jamais de double marque (\eng{more better}), jamais \eng{two times bigger than} à l'écrit.
\item Quantité : \eng{as much as} (indénombrable) / \eng{as many as} (dénombrable).
\item \eng{actually} = \emph{en fait} ; \emph{actuellement} = \eng{currently}.
\end{itemize}
\end{aretenir}

\coursec{Entraînement guidé et production}

Nous savons désormais distinguer les structures de comparaison anglaises et repérer les pièges que le français nous tend (\eng{more and more}, \eng{twice as much}, l'ordre des mots dans \eng{the more\ldots the more}). Il est temps de passer à la pratique. Cette dernière section vous entraîne, pas à pas, à mobiliser les quatre structures de la leçon dans des exercices de difficulté croissante, jusqu'à la production d'un paragraphe complet — exactement ce que l'on attend de vous au baccalauréat.

\courssub{La méthode générale : du mot déclencheur à la structure}

Avant de vous lancer dans les exercices, retenez la démarche qui garantit la réussite dans tous les exercices de transformation et de traduction.

\begin{methode}[Réussir un exercice de transformation ou de traduction]
\begin{enumerate}
\item \textbf{Lire} la phrase française (ou anglaise) en entier et \textbf{repérer le mot déclencheur} : \emph{aussi\ldots que}, \emph{moins\ldots que}, \emph{plus\ldots plus}, \emph{de plus en plus}, \emph{deux fois}, \emph{autant de}.
\item \textbf{Choisir la structure anglaise correspondante} à l'aide du tableau récapitulatif ci-dessous.
\item \textbf{Construire la phrase} en respectant l'ordre anglais : sujet + verbe + structure + complément.
\item \textbf{Vérifier} : adjectif court au comparatif avec \emph{-er} ou \emph{more} ? Pas de double marque (\eng{more better} est interdit) ? \emph{Than} et non \emph{that} ou \emph{then} ?
\item \textbf{Relire} la phrase à voix haute : sonne-t-elle naturellement anglaise ?
\end{enumerate}
\end{methode}

\begin{center}
\begin{tabularx}{\linewidth}{|X|X|X|}
\hline
\rowcolor{popblueL} Mot déclencheur français & Structure anglaise & Exemple \\
\hline
aussi\ldots que & as + adj + as & as rich as \\
\hline
moins\ldots que & less + adj + than / not as\ldots as & less crowded than \\
\hline
plus\ldots plus & the + comp., the + comp. & the more\ldots the more \\
\hline
de plus en plus & more and more / comp. + and + comp. & more and more expensive \\
\hline
deux fois, trois fois & twice / three times as\ldots as & twice as big as \\
\hline
autant de\ldots que & as much / as many\ldots as & as much money as \\
\hline
\end{tabularx}
\end{center}

\begin{popfigure}
\begin{tikzpicture}[
  node distance=0.9cm and 1.1cm,
  question/.style={rectangle, rounded corners, draw=popdark, fill=popblueL, text width=4.2cm, align=center, font=\bfseries},
  idea/.style={rectangle, rounded corners, draw=popblue, fill=popcream, text width=3.6cm, align=center, font=\small},
  struct/.style={rectangle, draw=poporange, fill=poporangeL, text width=3.4cm, align=center, font=\small\bfseries},
  fleche/.style={-{Stealth[length=3mm]}, thick, popdark}
]
\node[question] (q) {Que veux-tu exprimer ?};
\node[idea, below left=1.1cm and 2.2cm of q] (eg) {une égalité};
\node[idea, below left=1.1cm and 0.2cm of q] (prop) {une proportion};
\node[idea, below right=1.1cm and 0.2cm of q] (evol) {une évolution};
\node[idea, below right=1.1cm and 2.2cm of q] (mult) {un multiple};
\node[struct, below=0.8cm of eg] (s1) {as\ldots as};
\node[struct, below=0.8cm of prop] (s2) {the more\ldots the more};
\node[struct, below=0.8cm of evol] (s3) {more and more};
\node[struct, below=0.8cm of mult] (s4) {twice as\ldots as};
\draw[fleche] (q) -- (eg);
\draw[fleche] (q) -- (prop);
\draw[fleche] (q) -- (evol);
\draw[fleche] (q) -- (mult);
\draw[fleche, poporange] (eg) -- (s1);
\draw[fleche, poporange] (prop) -- (s2);
\draw[fleche, poporange] (evol) -- (s3);
\draw[fleche, poporange] (mult) -- (s4);
\end{tikzpicture}
\legende{Organigramme de choix : de l'idée française à la structure anglaise.}
\end{popfigure}

\courssub{Exercice 1 : compléter avec as\ldots as, not as\ldots as, less\ldots than}

Commençons par un exercice de manipulation contrôlée. Rappelez-vous : \eng{as\ldots as} exprime l'égalité, \eng{not as\ldots as} et \eng{less\ldots than} expriment l'infériorité.

\begin{exoresolu}[Exercice 1 — Compléter]
Complétez les phrases avec \eng{as\ldots as}, \eng{not as\ldots as} ou \eng{less\ldots than}.

1. A taxi is \ldots\ldots fast \ldots\ldots a bendskin in heavy traffic. (égalité)

2. The north of Cameroon is \ldots\ldots populated \ldots\ldots the south. (infériorité)

3. My new phone is \ldots\ldots expensive \ldots\ldots my old one. (infériorité)

4. English is \ldots\ldots important \ldots\ldots French in Cameroon. (égalité)
\tcblower
\textbf{Solutions détaillées :}

1. \eng{A taxi is as fast as a bendskin in heavy traffic.} — L'égalité exige \eng{as + adjectif + as}, sans modification de l'adjectif.

2. \eng{The north of Cameroon is not as populated as the south.} — On peut aussi dire \eng{less populated than} ; les deux sont corrects avec un adjectif long.

3. \eng{My new phone is less expensive than my old one.} — Avec un adjectif long (\emph{expensive}), \eng{less\ldots than} est la forme la plus naturelle ; \eng{not as expensive as} reste possible.

4. \eng{English is as important as French in Cameroon.} — Égalité : adjectif à la forme positive entre les deux \emph{as}.
\end{exoresolu}

\begin{attention}[Piège : ne jamais mettre le comparatif entre les deux as]
Les francophones écrivent souvent \eng{as better as} ou \eng{as more expensive as} par contamination du français « aussi meilleur que ». C'est faux : entre \emph{as\ldots as}, l'adjectif reste à la \textbf{forme positive} : \eng{as good as}, \eng{as expensive as}. De même, n'écrivez jamais \eng{as\ldots than} : le deuxième terme de l'égalité est toujours \emph{as}, jamais \emph{than}.
\end{attention}

\courssub{Exercice 2 : traduction guidée du français vers l'anglais}

La traduction est l'épreuve reine : elle oblige à choisir la bonne structure et à respecter l'ordre anglais.

\begin{exercice}[Exercice 2 — Traduire en anglais]
Traduisez les phrases suivantes en utilisant la structure indiquée entre parenthèses.

1. Douala est aussi chaude que Garoua. (\eng{as\ldots as})

2. Un enseignant gagne moins qu'un banquier. (\eng{less\ldots than})

3. Plus tu travailles, plus tu réussis. (\eng{the more\ldots the more})

4. La vie devient de plus en plus chère à Yaoundé. (\eng{more and more})

5. Ma sœur est deux fois plus âgée que mon frère. (\eng{twice as\ldots as})
\end{exercice}

\begin{corrige}[Exercice 2]
1. \eng{Douala is as hot as Garoua.}

2. \eng{A teacher earns less than a banker.} — Ici \eng{less than} suffit, car on compare des montants (verbe \emph{earn} + adverbe de quantité).

3. \eng{The more you work, the more you succeed.} ou, plus naturellement, \eng{The harder you work, the more successful you become.}

4. \eng{Life is becoming more and more expensive in Yaoundé.} — Adjectif long : \eng{more and more + adjectif}.

5. \eng{My sister is twice as old as my brother.} — Attention : « deux fois plus âgée » se dit \eng{twice as old as}, jamais \eng{twice older than} dans un anglais soigné.
\end{corrige}

\courssub{Exercice 3 : réécrire avec the more\ldots the more}

Cette structure de proportion est la préférée des examinateurs. Observons d'abord quelques exemples économiques, dans l'esprit du chapitre.

\begin{exemplebox}[La proportion dans la vie économique]
\eng{The more goods a trader sells, the more profit he makes.} « Plus un commerçant vend de marchandises, plus il fait de bénéfices. »

\eng{The higher the price of fuel, the more expensive transport becomes.} « Plus le prix du carburant est élevé, plus le transport devient cher. »

\eng{The harder farmers work, the better the harvest.} « Plus les agriculteurs travaillent dur, meilleure est la récolte. » — Remarquez : quand le verbe est \emph{be}, on peut l'omettre dans la deuxième partie (\eng{the better the harvest} au lieu de \eng{the better the harvest is}).
\end{exemplebox}

\begin{exercice}[Exercice 3 — Transformer]
Réécrivez chaque paire d'idées en une seule phrase avec \eng{the\ldots the}.

1. You save money early. You become rich quickly.

2. A student reads many books. He writes well.

3. The city grows. Housing becomes scarce.
\end{exercice}

\begin{corrige}[Exercice 3]
1. \eng{The earlier you save money, the richer you become.} — Adverbe court \emph{early} : comparatif \emph{earlier}.

2. \eng{The more books a student reads, the better he writes.} — \emph{Well} a un comparatif irrégulier : \emph{better}.

3. \eng{The more the city grows, the scarcer housing becomes.} — \emph{Scarce} prend \emph{-er} : \emph{scarcer}.
\end{corrige}

\courssub{Exercice 4 : corriger les erreurs typiques des francophones}

Voici l'exercice le plus formateur : chaque phrase contient exactement une erreur commise très souvent par les élèves francophones.

\begin{exoresolu}[Exercice 4 — Trouver et corriger l'erreur]
1. \eng{Yaoundé is as bigger as Douala.}

2. \eng{The more you study, the more better marks you get.}

3. \eng{Petrol is more and more and more expensive.}

4. \eng{This market is twice bigger than the old one.}

5. \eng{She earns less money that her brother.}
\tcblower
\textbf{Corrections :}

1. \eng{Yaoundé is as big as Douala.} — Entre \emph{as\ldots as}, forme positive, pas de comparatif.

2. \eng{The more you study, the better your marks.} — Double comparatif interdit : \eng{more better} est une faute classique ; \emph{better} est déjà le comparatif de \emph{good}.

3. \eng{Petrol is more and more expensive.} — Un seul \eng{more and more} ; la répétition triple est une contamination du français « de plus en plus » mal décomposé.

4. \eng{This market is twice as big as the old one.} — Le multiple exige \eng{twice as\ldots as}, pas \eng{twice bigger than}.

5. \eng{She earns less money than her brother.} — \emph{Than} avec \emph{th-} ; \emph{that} est un démonstratif. Prononciation : « zane » faible, jamais « zate ».
\end{exoresolu}

\begin{savaistu}[La transformation, un classique du Bac]
Au baccalauréat, les questions de grammaire portent très souvent sur la transformation de phrases. Savoir passer de \eng{X is bigger than Y} à \eng{Y is not as big as X} est un classique absolu : les deux phrases ont le même sens, mais la structure change. Entraînez-vous dans les deux sens : supériorité vers égalité négative, et égalité vers infériorité. Par exemple : \eng{A car is more expensive than a motorcycle.} devient \eng{A motorcycle is not as expensive as a car.}
\end{savaistu}

\courssub{Exercice 5 : comparer deux villes camerounaises}

Passons à la production semi-guidée. Lisez d'abord le texte modèle ci-dessous : il compare Buea et Bamenda en utilisant les quatre structures de la leçon.

\begin{textebox}[Two towns, two rhythms]
Buea and Bamenda are both important towns in the English-speaking regions of Cameroon, but they offer very different lives. Buea is not as large as Bamenda, yet it is just as lively, thanks to its university and its many students. The climate of Buea is less hot than the climate of Douala, because the town lies on the slopes of Mount Cameroon. Bamenda, on the other hand, is a major commercial centre: the more traders arrive from the North-West villages, the busier its markets become. Food there is twice as cheap as in Yaoundé, according to many travellers. However, life in both towns is becoming more and more expensive, and transport fares are getting higher and higher. Young graduates say that the more qualifications they have, the easier it is to find a job in these growing towns. Choosing between Buea and Bamenda is therefore not as simple as it seems: each town has its own rhythm, its own opportunities and its own charm.

\textbf{Glossaire :} \emph{slopes} (n., pentes) ; \emph{traders} (n., commerçants) ; \emph{fares} (n., tarifs de transport) ; \emph{qualifications} (n., diplômes) ; \emph{charm} (n., charme).
\end{textebox}

\begin{exercice}[Exercice 5 — À votre tour]
Comparez deux villes camerounaises de votre choix (par exemple Douala et Yaoundé, ou Garoua et Maroua) en cinq ou six phrases. Utilisez \textbf{au moins trois structures différentes} parmi : \eng{as\ldots as}, \eng{not as\ldots as}, \eng{less\ldots than}, \eng{the more\ldots the more}, \eng{more and more}, \eng{twice as\ldots as}. Soulignez chaque structure utilisée.
\end{exercice}

\begin{corrige}[Exercice 5 — Production possible]
\eng{Douala is larger than Yaoundé, but Yaoundé is not as noisy as the economic capital. Life in Douala is less calm than life in Yaoundé because of the port and the traffic. The more businesses open in Douala, the more workers arrive from the villages. Rent in Bonanjo can be twice as expensive as in some quarters of Yaoundé. Both cities are becoming more and more crowded, and finding a house is getting harder and harder.}
\end{corrige}

\courssub{Production finale : comparer deux métiers}

Terminons par la tâche intégratrice, directement liée au chapitre \emph{Economic Life and Occupations}.

\begin{attention}[Varier les structures : un critère d'excellence]
Dans la grille d'évaluation de la production écrite, la \textbf{variété des structures} est un critère explicite. Un paragraphe qui répète cinq fois \eng{more\ldots than} obtient une note moyenne ; un paragraphe qui alterne égalité, proportion, gradation et multiple démontre la maîtrise de la langue. Variez donc délibérément vos comparaisons.
\end{attention}

\begin{experience}[Production écrite guidée : Two jobs compared]
Rédigez un paragraphe de huit à dix phrases comparant deux métiers (par exemple \emph{a teacher} et \emph{a banker}, ou \emph{a farmer} et \emph{a nurse}) selon trois critères : le salaire, la difficulté, les horaires.

\textbf{Contraintes :} utilisez les quatre structures de la leçon, au moins une fois chacune : \eng{as\ldots as} ou \eng{not as\ldots as} ; \eng{less\ldots than} ; \eng{the more\ldots the more} ; \eng{more and more} ou \eng{twice as\ldots as}.

\textbf{Plan suggéré :} 1) présenter les deux métiers ; 2) comparer les salaires ; 3) comparer la difficulté ; 4) comparer les horaires ; 5) conclure par une préférence justifiée.
\end{experience}

\begin{exemplebox}[Modèle de production]
\eng{A teacher and a banker both play important roles in our society, but their working lives are quite different. A banker earns more than a teacher; in fact, his salary can be twice as high as a teacher's salary. However, teaching is not as easy as many people think: the more students a teacher has in class, the harder his work becomes. A banker's job is less tiring physically, but it is more and more stressful because of targets and long hours at the office. Teachers enjoy longer holidays, so their free time is not as limited as that of bankers. As for me, the more I think about the two professions, the more I admire teachers, because shaping young minds is as valuable as earning a big salary.}
\end{exemplebox}

\begin{aretenir}[L'essentiel de la section]
\begin{itemize}
\item Méthode : repérer le \cle{mot déclencheur} français (\emph{aussi}, \emph{moins}, \emph{plus\ldots plus}, \emph{de plus en plus}, \emph{fois}), choisir la structure, construire, vérifier.
\item \eng{As\ldots as} : adjectif à la forme positive ; jamais de comparatif entre les deux \emph{as}.
\item \eng{The more\ldots the more} : deux comparatifs, jamais de double marque (\eng{more better} est interdit).
\item \eng{Twice as\ldots as} pour les multiples ; \eng{more and more} pour la gradation.
\item À l'écrit, \cle{varier les structures} de comparaison est un critère d'excellence au baccalauréat.
\end{itemize}
\end{aretenir}

\newpage

\coursec{Activité d'intégration}

\courssub{Objectifs de l'activité}

À l'issue de cette activité d'intégration, tu seras capable de :
\begin{itemize}
\item mobiliser dans une production authentique les \cle{cinq formes de comparaison} étudiées dans la leçon : l'égalité (\eng{as... as}), l'inégalité (\eng{not as... as}, \eng{less... than}), la proportion (\eng{the more... the more}), la gradation (\eng{more and more}) et la multiplication (\eng{twice as... as}) ;
\item comparer deux métiers exercés au Cameroun en utilisant un lexique économique précis (salaire, revenu, formation, conditions de travail) ;
\item rédiger un paragraphe comparatif cohérent de 8 à 10 lignes et le présenter à l'oral avec une prononciation claire ;
\item repérer et nommer les structures comparatives dans la production d'autrui (compétence d'écoute active).
\end{itemize}

\courssub{Situation}

Ton établissement organise sa \emph{Career Guidance Week} (semaine d'orientation professionnelle). Le club d'anglais, dont tu es membre, a été chargé de publier dans le journal de l'école une série de courts articles comparant des métiers que les jeunes Camerounais envisagent d'exercer. Le thème de cette semaine : \emph{“Two jobs, two lives: which one for you?”}. Pour préparer les articles, le professeur d'anglais distribue le texte d'amorce suivant, inspiré d'une enquête menée auprès d'élèves de Terminale à Yaoundé et à Bamenda.

\begin{textebox}[Two jobs, two lives — reading input]
In Cameroon today, choosing an occupation is not as simple as it used to be. Take the job of a secondary school teacher and that of a market trader in Mokolo market, Yaoundé. A teacher's monthly salary is regular, but it is not as high as many people imagine. A successful trader can earn twice as much money as a teacher, especially during the festive season. However, the more customers a trader has, the longer she works, and her income is never guaranteed: some weeks are good, others are not.

Teaching is becoming more and more demanding: classes are getting larger and larger, and the more experienced a teacher becomes, the more responsibilities he is given. Yet many young graduates still prefer it because it offers stability. Trading, on the other hand, is less secure than teaching, but it can be far more profitable. In the end, the choice depends on one's personality: are you as patient as a teacher must be, or as bold as a trader needs to be?
\end{textebox}

\textbf{Glossaire :} \emph{salary} (n.m.) : salaire ; \emph{income} (n.m.) : revenu ; \emph{guaranteed} (adj.) : garanti ; \emph{demanding} (adj.) : exigeant ; \emph{bold} (adj.) : audacieux ; \emph{profitable} (adj.) : rentable ; \emph{festive season} (loc. nom.) : période des fêtes.

\courssub{Consignes}

Travail en groupes de quatre. Durée totale : 50 minutes (35 minutes de préparation écrite, 15 minutes de présentations orales).

\begin{enumerate}
\item \textbf{Compréhension du texte d'amorce (5 min).} Lis le texte \emph{“Two jobs, two lives”}. Souligne toutes les structures comparatives et classe-les dans un tableau à cinq colonnes : égalité, inégalité, proportion, gradation, multiplication. Vérifie que chaque catégorie est représentée.
\item \textbf{Choix des métiers (3 min).} Avec ton groupe, choisissez deux métiers exercés au Cameroun : \eng{teacher vs trader}, \eng{nurse vs taxi driver}, \eng{farmer vs bank clerk}, \eng{tailor vs civil servant}, ou tout autre couple pertinent. Les deux métiers doivent être suffisamment différents pour être comparés.
\item \textbf{Remue-méninges lexical (5 min).} Complétez ensemble une fiche de vocabulaire : pour chaque métier, notez au moins trois éléments comparables (salaire, horaires, formation, prestige, sécurité de l'emploi). Utilisez des adjectifs graduables : \eng{high, demanding, secure, profitable, stressful, respected}.
\item \textbf{Rédaction du paragraphe (15 min).} Rédigez un paragraphe comparatif de 8 à 10 lignes intitulé \emph{“Job A vs Job B”}. Il doit contenir \textbf{obligatoirement} :
\begin{itemize}
\item une comparaison d'\cle{égalité} : \eng{as... as} ;
\item une comparaison d'\cle{inégalité} : \eng{not as... as} ou \eng{less... than} ;
\item une \cle{proportion} : \eng{the more... the more...} ;
\item une \cle{gradation} : \eng{more and more} ou un comparatif doublé (\eng{larger and larger}) ;
\item une \cle{multiplication} : \eng{twice as... as} (ou \eng{three times as... as}).
\end{itemize}
Chaque structure doit porter sur un aspect différent du métier (ne comparez pas cinq fois le salaire).
\item \textbf{Relecture collective (5 min).} Échangez votre texte avec le groupe voisin. Vérifiez : la forme des structures (deux \eng{as} bien présents, \eng{the} devant chaque comparatif de proportion), l'accord des verbes, l'orthographe des comparatifs (\eng{bigger}, \eng{more secure}). Corrigez les erreurs relevées.
\item \textbf{Présentation orale (15 min).} Deux membres du groupe lisent le texte à voix haute, avec une prononciation soignée (attention à \eng{the more... the more} : « dé mor... dé mor »). Pendant chaque présentation, les autres élèves cochent sur une grille les cinq structures dès qu'ils les entendent et les citent ensuite à l'oral.
\item \textbf{Vote et affichage (2 min).} La classe vote pour le texte le plus riche et le plus correct. Le texte gagnant est affiché au mur de la classe dans la rubrique \emph{English Corner}.
\end{enumerate}

\courssub{Ressources à mobiliser}

\begin{aretenir}[Les cinq structures à utiliser]
\begin{center}
\begin{tabularx}{\linewidth}{|l|X|X|}
\hline
\rowcolor{popblueL} \textbf{Forme} & \textbf{Structure} & \textbf{Exemple} \\
\hline
Égalité & as + adj. + as & A nurse is as busy as a doctor. \\
\hline
Inégalité & not as + adj. + as / less + adj. + than & Trading is less secure than teaching. \\
\hline
Proportion & the + comp. + sujet + verbe, the + comp. + sujet + verbe & The more you sell, the more you earn. \\
\hline
Gradation & more and more + adj. / comp. + and + comp. & Farming is becoming more and more modern. \\
\hline
Multiplication & twice / three times + as + adj. + as & A trader can earn twice as much as a clerk. \\
\hline
\end{tabularx}
\end{center}
\end{aretenir}

\textbf{Lexique utile (economic life and occupations) :} \eng{salary, wages, income, profit, to earn, to employ, to hire, working hours, job security, training, qualification, demanding, rewarding, stressful, self-employed, civil servant}.

\begin{attention}[Pièges à éviter dans ta production]
\begin{itemize}
\item Ne dis jamais \eng{more and more bigger} : on choisit soit \eng{bigger and bigger} (adjectif court), soit \eng{more and more + adjectif long} (\eng{more and more expensive}).
\item Dans la proportion, les deux membres prennent \eng{the} : \eng{the more..., the more...} — jamais \eng{more you work, more you earn}.
\item \eng{Twice as much as} (et non \eng{two times more than}, calqué du français « deux fois plus que »). De même : \eng{three times as many passengers as}, et non \eng{three times more passengers than}.
\item Avec \eng{twice as... as}, choisis le bon quantifieur : \eng{twice as \textbf{much} money as} (indénombrable) mais \eng{twice as \textbf{many} customers as} (dénombrable).
\item Pas de double comparaison : \eng{more better} est incorrect ; dis \eng{better}. De même, \eng{more easier} est fautif : dis \eng{easier}.
\item Dans \eng{as... as}, l'adjectif reste à la forme positive : \eng{as good as}, jamais \eng{as better as}.
\end{itemize}
\end{attention}

\courssub{Proposition de production et corrigé commenté}

\begin{exoresolu}[Modèle de paragraphe comparatif : “Nurse vs Taxi Driver”]
\textbf{Énoncé :} rédige le paragraphe comparatif de ton groupe en respectant les cinq structures obligatoires.

\tcblower

\textbf{Production modèle :}

\emph{“Nurse vs Taxi Driver. In Douala, a nurse and a taxi driver both work hard, but their lives are quite different. A nurse's job is as demanding as a doctor's, because she cares for patients day and night. However, her salary is not as high as many people think, and a good taxi driver in Bonabéri can sometimes earn twice as much money as a nurse in a public hospital. The more passengers a driver carries, the more he earns, but his income depends on the season and the price of fuel. Nursing, on the other hand, offers more security: it is becoming more and more respected, and hospitals need more and more qualified staff. In the end, the choice is yours: do you prefer freedom on the road or stability in a uniform?”}

\textbf{Commentaire des choix de langue :}
\begin{itemize}
\item \textbf{Égalité} : \eng{as demanding as a doctor's} — adjectif à la forme positive entre les deux \eng{as} ; le génitif \eng{a doctor's} évite de répéter \eng{job}.
\item \textbf{Inégalité} : \eng{not as high as many people think} — variante naturelle de \eng{less high than}, plus idiomatique en anglais courant.
\item \textbf{Multiplication} : \eng{twice as much money as} — \eng{much} car \eng{money} est indénombrable ; structure correcte \eng{twice as + quantifieur + nom + as}.
\item \textbf{Proportion} : \eng{the more passengers a driver carries, the more he earns} — \eng{the} devant chaque comparatif ; l'ordre sujet--verbe est conservé dans les deux membres, sans inversion.
\item \textbf{Gradation} : \eng{more and more respected} (adjectif long) et \eng{more and more qualified staff} — deux occurrences qui renforcent l'idée d'évolution progressive.
\item \textbf{Cohérence} : chaque structure porte sur un aspect différent (exigence, salaire, revenu variable, prestige), et la conclusion interrogative relance le débat, ce qui convient au genre « article d'orientation ».
\end{itemize}
\end{exoresolu}

\courssub{Bilan de l'activité}

Cette activité t'a permis de passer de la \cle{reconnaissance} des structures comparatives à leur \cle{production} en contexte réel. Retiens les points essentiels :
\begin{itemize}
\item les cinq formes de comparaison (égalité, inégalité, proportion, gradation, multiplication) ont chacune une \textbf{forme fixe} qu'il faut respecter à la lettre : deux \eng{as}, deux \eng{the}, pas de double comparatif ;
\item dans un texte comparatif authentique, on \textbf{varie les structures} et les aspects comparés pour éviter la monotonie ;
\item à l'oral, la proportion \eng{the more... the more} se prononce de façon fluide et rythmée — entraîne-toi à la dire sans hésitation ;
\item ces structures te seront utiles à l'examen, tant en compréhension de l'écrit (repérage dans un texte) qu'en expression écrite (essai, article, lettre).
\end{itemize}

\begin{experience}[Pour aller plus loin : mini-débat]
En prolongement, organise un débat de cinq minutes : \emph{“Which job is better for a young Cameroonian today?”} Chaque groupe défend le métier de son choix en utilisant au moins trois structures comparatives différentes. L'auditoire note chaque structure correctement employée : une structure correcte rapporte un point à l'équipe.
\end{experience}

\newpage

\coursec{Méthodes et exercices résolus}

\courssub{Savoir-faire 1 : Identifier et construire la comparaison d'égalité et d'inégalité}

\begin{propriete}[Les structures de base]
\begin{itemize}
\item \textbf{Égalité} : \cle{as + adjectif/adverbe + as} : \eng{This market is as busy as the one in Bafoussam.} « Ce marché est aussi animé que celui de Bafoussam. »
\item \textbf{Infériorité} : \cle{not as/so + adjectif/adverbe + as} : \eng{The village shop is not as busy as the town market.} « La boutique du village n'est pas aussi animée que le marché de la ville. »
\item \textbf{Avec un nom indénombrable} : \cle{as much + nom + as} : \eng{She earns as much money as her brother.} « Elle gagne autant d'argent que son frère. »
\item \textbf{Avec un nom dénombrable} : \cle{as many + nom pluriel + as} : \eng{He employs as many workers as his competitor.} « Il emploie autant d'ouvriers que son concurrent. »
\end{itemize}
\end{propriete}

\begin{methode}[Bien construire « as... as » et « not as/so... as »]
\begin{enumerate}
\item Repérer d'abord l'\textbf{adjectif} ou l'\textbf{adverbe} à comparer : c'est lui qui se place entre \cle{as} et \cle{as}.
\item Appliquer la structure d'égalité : \eng{as + adjectif/adverbe + as + terme de comparaison}. Exemple : \eng{Trade in Douala is as active as trade in Lagos.} « Le commerce à Douala est aussi actif qu'à Lagos. »
\item Pour la négation (infériorité), utiliser \eng{not as/so ... as} : \eng{The second shop is not as busy as the first one.} « La deuxième boutique n'est pas aussi animée que la première. »
\item Avec un nom, choisir \cle{as much + nom indénombrable + as} ou \cle{as many + nom dénombrable + as} : \eng{She earns as much money as her brother.} / \eng{He employs as many workers as his competitor.}
\item Avec un nom singulier précédé d'un adjectif, l'ordre est \eng{as + adjectif + a + nom + as} : \eng{It is as difficult a job as mining.} « C'est un métier aussi difficile que celui de mineur. »
\item Ne jamais mettre de comparatif (\eng{-er} ou \eng{more}) entre les deux \eng{as} : on dit \eng{as cheap as}, jamais \eng{as cheaper as}.
\end{enumerate}
\end{methode}

\begin{exoresolu}[Equality and inequality in context]
\textbf{Énoncé.} Complete the sentences with the correct structure (\eng{as...as}, \eng{not as/so...as}, \eng{as much...as}, \eng{as many...as}).

a) A taxi driver in Yaoundé works .................... hours .................... a bendskin rider. (long)

b) My uncle's cocoa farm does not produce .................... cocoa .................... my father's. (much)

c) The new supermarket sells .................... goods .................... the old market. (many)

d) This job is .................... tiring .................... working in a bank. (not / exhausting)

\tcblower
\textbf{Solution détaillée.}

a) Il s'agit d'une égalité portant sur un adverbe (\eng{long}). Structure : \eng{as + adverbe + as}. Réponse : \eng{A taxi driver in Yaoundé works as long hours as a bendskin rider.} « Un chauffeur de taxi à Yaoundé travaille de longues heures, tout comme un moto-taxi. »

b) \eng{Cocoa} est indénombrable et la phrase est négative : on utilise \eng{not as much + nom + as}. Réponse : \eng{My uncle's cocoa farm does not produce as much cocoa as my father's.} « La plantation de cacao de mon oncle ne produit pas autant de cacao que celle de mon père. »

c) \eng{Goods} est dénombrable pluriel : \eng{as many + nom pluriel + as}. Réponse : \eng{The new supermarket sells as many goods as the old market.} « Le nouveau supermarché vend autant de marchandises que le vieux marché. »

d) Infériorité avec adjectif long : \eng{not as + adjectif + as}. Réponse : \eng{This job is not as exhausting as working in a bank.} « Ce travail n'est pas aussi épuisant que travailler dans une banque. »

\textbf{Conclusion.} Le choix entre \eng{much} et \eng{many} dépend uniquement de la nature du nom (indénombrable ou dénombrable), et l'adjectif reste toujours à la forme simple entre les deux \eng{as}.
\end{exoresolu}

\courssub{Savoir-faire 2 : Employer la proportion « the more... the more »}

\begin{propriete}[La double proportion]
La structure \cle{the + comparatif + sujet + verbe, the + comparatif + sujet + verbe} exprime deux phénomènes qui varient ensemble : \eng{The harder you work, the more you earn.} « Plus tu travailles dur, plus tu gagnes. » La première proposition prend le \textbf{présent} (jamais le futur) ; la seconde peut prendre le futur : \eng{The more you study, the more you will succeed.}
\end{propriete}

\begin{methode}[Construire la double proportion]
\begin{enumerate}
\item Identifier les deux éléments qui varient ensemble (cause et conséquence).
\item Appliquer le schéma : \cle{the + comparatif + sujet + verbe, the + comparatif + sujet + verbe}. Exemple : \eng{The harder you work, the more you earn.} « Plus tu travailles dur, plus tu gagnes. »
\item Pour les adjectifs courts, le comparatif est en \eng{-er} : \eng{The sooner, the better.} « Le plus tôt sera le mieux. »
\item Pour les adjectifs longs, utiliser \eng{more} : \eng{The more expensive the goods are, the fewer customers the shop attracts.} « Plus les marchandises sont chères, moins la boutique attire de clients. »
\item Attention à l'ordre des mots : le comparatif vient \textbf{en tête} de chaque proposition, puis sujet + verbe (pas d'inversion).
\item Ne jamais traduire mot à mot le français « plus... plus... » par \eng{more... more...} seul : il faut toujours \textbf{the} devant chaque comparatif.
\end{enumerate}
\end{methode}

\begin{exoresolu}[Expressing proportion]
\textbf{Énoncé.} Rewrite each pair of ideas using the \eng{the... the...} structure.

a) You save money early. You become independent early.

b) A product is cheap. Many people buy it.

c) Translate into English : « Plus un commerçant est honnête, plus il a de clients fidèles. »

\tcblower
\textbf{Solution détaillée.}

a) L'adverbe \eng{early} a pour comparatif \eng{earlier}. Réponse : \eng{The earlier you save money, the earlier you become independent.} « Plus tôt tu économises, plus tôt tu deviens indépendant. » On place le comparatif en tête, suivi du sujet puis du verbe, sans inversion.

b) \eng{Cheap} est un adjectif court : \eng{cheaper}. Pour la deuxième proposition, on dit \eng{the more people buy it}. Réponse : \eng{The cheaper a product is, the more people buy it.} « Moins un produit est cher, plus les gens l'achètent. »

c) « Honnête » se dit \eng{honest} (comparatif \eng{more honest}) ; « clients fidèles » = \eng{loyal customers}. Réponse : \eng{The more honest a trader is, the more loyal customers he has.} On remarque que le français répète « plus » seul, alors que l'anglais exige \textbf{the + comparatif} dans les deux propositions.

\textbf{Conclusion.} La structure est symétrique : deux propositions, deux fois \eng{the + comparatif}, et jamais d'inversion du sujet.
\end{exoresolu}

\courssub{Savoir-faire 3 : Exprimer la gradation et la multiplication}

\begin{propriete}[Gradation et multiples]
\begin{itemize}
\item \textbf{De plus en plus} : comparatif répété avec \cle{and} (\eng{higher and higher}) ou \cle{more and more + adjectif long}.
\item \textbf{De moins en moins} : \cle{less and less + adjectif/nom indénombrable}, \cle{fewer and fewer + nom dénombrable}.
\item \textbf{Fois plus} : \cle{twice / three times + as much/many ... as} ; \textbf{le double de} : \cle{double the + nom} ; \textbf{la moitié de} : \cle{half as much/many ... as}.
\end{itemize}
\end{propriete}

\begin{methode}[Gradation croissante et multiples]
\begin{enumerate}
\item \textbf{Gradation croissante} : répéter le comparatif avec \cle{and} : \eng{Prices are getting higher and higher.} « Les prix deviennent de plus en plus élevés. » Pour les adjectifs longs : \eng{more and more + adjectif} : \eng{Life in Douala is becoming more and more expensive.} « La vie à Douala devient de plus en plus chère. »
\item \textbf{Gradation décroissante} : \eng{less and less + adjectif} ou \eng{fewer and fewer + nom dénombrable} : \eng{Fewer and fewer young people want to become farmers.} « De moins en moins de jeunes veulent devenir agriculteurs. »
\item \textbf{Multiplication} : \cle{twice / three times / half} + \eng{as much/many ... as} : \eng{This shop earns twice as much as the other one.} « Cette boutique gagne deux fois plus que l'autre. »
\item Avec un nom : \eng{double / twice / three times + the + nom} : \eng{He pays double the price.} « Il paie le double du prix. »
\item Réflexe : vérifier si le nom est dénombrable (\eng{twice as many workers}) ou indénombrable (\eng{twice as much money}).
\end{enumerate}
\end{methode}

\begin{exoresolu}[Gradation and multiples]
\textbf{Énoncé.} a) Rewrite : « The cost of transport increases continuously. » using \eng{higher and higher}.

b) Complete : The new factory produces .................... (three times / many) shoes .................... the old workshop.

c) Translate : « Le loyer à Bamenda coûte deux fois moins cher qu'à Douala. »

\tcblower
\textbf{Solution détaillée.}

a) « De plus en plus » avec un adjectif court se rend par la répétition du comparatif. Réponse : \eng{The cost of transport is getting higher and higher.} « Le coût du transport devient de plus en plus élevé. » Le verbe \eng{get} + comparatif exprime le changement progressif.

b) \eng{Shoes} est dénombrable : il faut \eng{many}, et la structure est \eng{three times as many ... as}. Réponse : \eng{The new factory produces three times as many shoes as the old workshop.} « La nouvelle usine produit trois fois plus de chaussures que l'ancien atelier. »

c) « Deux fois moins cher » se rend par \eng{half as much as} ou \eng{half the price}. Réponse : \eng{Rent in Bamenda costs half as much as in Douala.} ou \eng{Rent in Bamenda is half the price of rent in Douala.} « Le loyer à Bamenda coûte moitié moins cher qu'à Douala. » On évite le calque \eng{two times less expensive}, qui est incorrect en anglais standard.

\textbf{Conclusion.} Pour la multiplication, l'anglais raisonne en « fois autant que » (\eng{twice as much as}) et en « moitié » (\eng{half as much as}), jamais en « fois moins ».
\end{exoresolu}

\courssub{Savoir-faire 4 : Éviter les calques du français}

\begin{center}
\begin{tabularx}{\linewidth}{|X|X|X|}
\hline
\rowcolor{popblueL} Français & English & Remarque \\
\hline
autant d'argent & as much money & indénombrable \\
\hline
autant de clients & as many customers & dénombrable \\
\hline
de plus en plus cher & more and more expensive & adjectif long \\
\hline
de moins en moins de temps & less and less time & indénombrable \\
\hline
le double du prix & double the price & jamais \eng{the double of} \\
\hline
deux fois plus grand que & twice as big as & forme simple après \eng{as} \\
\hline
plus... plus... & the more... the more... & \eng{the} obligatoire \\
\hline
\end{tabularx}
\end{center}

\begin{methode}[Traduire les comparaisons sans calquer le français]
\begin{enumerate}
\item « Autant de » : \eng{as much} (indénombrable) ou \eng{as many} (dénombrable), jamais \eng{as much of}.
\item « De plus en plus de » : \eng{more and more} ; « de moins en moins de » : \eng{less and less} (indénombrable) / \eng{fewer and fewer} (dénombrable).
\item « Le double de » : \eng{double the...} ou \eng{twice as much/many... as}, jamais \eng{the double of}.
\item « Deux fois plus grand que » : \eng{twice as big as} (et non \eng{two times bigger than} en anglais soigné).
\item « Plus... plus... » : \eng{the + comparatif... the + comparatif...}, avec \textbf{the} obligatoire.
\item « Comparé à » : \eng{compared to/with}, sans faux ami avec \eng{comparative}.
\end{enumerate}
\end{methode}

\begin{exoresolu}[Correct the mistakes]
\textbf{Énoncé.} Each sentence contains one mistake made by a French-speaking candidate. Correct it.

a) Yaoundé has the double of inhabitants of Buea.

b) The more you will study, the more you will succeed.

c) Petrol is becoming more and more cheaper.

d) My sister earns as much money than me.

\tcblower
\textbf{Solution détaillée.}

a) Calque de « le double de ». Correction : \eng{Yaoundé has double the number of inhabitants of Buea.} ou \eng{Yaoundé has twice as many inhabitants as Buea.} « Yaoundé a deux fois plus d'habitants que Buea. »

b) Après \eng{the more}, on n'emploie pas le futur \eng{will} dans la première proposition : le présent suffit. Correction : \eng{The more you study, the more you will succeed.} « Plus tu étudies, plus tu réussiras. »

c) Double marque du comparatif : \eng{more} + \eng{-er} est impossible. Correction : \eng{Petrol is becoming more and more expensive.} ou \eng{Petrol is becoming cheaper and cheaper.} « L'essence devient de plus en plus chère. »

d) Avec \eng{as...as}, le deuxième élément est \eng{as}, jamais \eng{than}. Correction : \eng{My sister earns as much money as me / as I do.} « Ma sœur gagne autant d'argent que moi. »

\textbf{Conclusion.} Chaque erreur vient d'un calque du français : identifier la structure anglaise complète avant d'écrire est le meilleur réflexe.
\end{exoresolu}

\courssub{Savoir-faire 5 : Utiliser ces structures dans un court texte}

\begin{methode}[Rédiger un paragraphe comparatif]
\begin{enumerate}
\item Choisir un sujet concret (deux villes, deux métiers, deux marchés) et noter deux ou trois points de comparaison.
\item Varier les structures : une égalité (\eng{as...as}), une inégalité (\eng{not as...as} ou comparatif), une proportion (\eng{the more... the more}) et un multiple (\eng{twice as... as}).
\item Relier les idées avec des connecteurs : \eng{however}, \eng{moreover}, \eng{as a result}, \eng{on the other hand}.
\item Relire en vérifiant : adjectif sans \eng{-er} entre \eng{as...as}, \eng{the} devant chaque comparatif de proportion, \eng{much/many} adapté au nom.
\item Conclure par une phrase d'appréciation personnelle.
\end{enumerate}
\end{methode}

\begin{exoresolu}[Guided composition]
\textbf{Énoncé.} Write a short paragraph (60–80 words) comparing the work of a market trader and a bank clerk in Cameroon, using at least three comparison structures seen in this lesson.

\tcblower
\textbf{Solution détaillée.}

\textbf{Plan.} 1) égalité sur les horaires ; 2) inégalité sur le revenu ; 3) proportion travail/gain ; 4) multiple sur les gains ; 5) gradation sur les exigences du métier.

\textbf{Paragraphe rédigé.}

\eng{A market trader in Douala works as many hours as a bank clerk, but their lives are quite different. The trader does not earn as regular a salary as the clerk, because his income depends on daily sales. However, the harder he works, the more customers he attracts, and a successful trader can earn twice as much money as a bank employee. On the other hand, the clerk's job is becoming more and more demanding, with less and less free time. In my opinion, both occupations are equally important for our economy.}

« Un commerçant au marché de Douala travaille autant d'heures qu'un employé de banque, mais leurs vies sont très différentes. Le commerçant ne gagne pas un salaire aussi régulier que l'employé, car son revenu dépend des ventes quotidiennes. Cependant, plus il travaille dur, plus il attire de clients, et un commerçant prospère peut gagner deux fois plus d'argent qu'un employé de banque. En revanche, le métier d'employé devient de plus en plus exigeant, avec de moins en moins de temps libre. À mon avis, les deux métiers sont tout aussi importants pour notre économie. »

\textbf{Vérification.} \eng{as many hours as} (dénombrable) ; \eng{as regular a salary as} (ordre : \eng{as + adjectif + a + nom}) ; \eng{the harder... the more...} (deux \eng{the}) ; \eng{twice as much money as} (indénombrable) ; \eng{more and more demanding}, \eng{less and less free time}. Toutes les structures sont correctes et variées.
\end{exoresolu}

\begin{aretenir}[L'essentiel de la leçon]
\begin{itemize}
\item Égalité : \eng{as + adjectif + as} ; inégalité : \eng{not as/so + adjectif + as}.
\item Noms : \eng{as much + indénombrable}, \eng{as many + dénombrable}.
\item Proportion : \eng{the + comparatif..., the + comparatif...} (présent dans la première proposition).
\item Gradation : \eng{higher and higher}, \eng{more and more expensive}, \eng{less and less}, \eng{fewer and fewer}.
\item Multiples : \eng{twice/three times as much/many ... as}, \eng{double the price}, \eng{half as much as}.
\end{itemize}
\end{aretenir}

\begin{attention}[Les 5 erreurs qui coûtent des points au Bac]
\begin{enumerate}
\item \textbf{Comparatif entre « as... as »} : écrire \eng{as cheaper as} au lieu de \eng{as cheap as}. L'adjectif reste à la forme simple.
\item \textbf{Oublier « the » dans la proportion} : \eng{More you work, more you earn} est faux ; il faut \eng{The more you work, the more you earn.}
\item \textbf{Calque de « le double de »} : \eng{the double of} n'existe pas ; dire \eng{double the price} ou \eng{twice as much as}.
\item \textbf{Confondre much et many} : \eng{as much workers} est incorrect ; \eng{workers} est dénombrable, donc \eng{as many workers}.
\item \textbf{Futur après « the more »} : \eng{The more you will study...} est une faute de temps ; la première proposition prend le présent : \eng{The more you study, the more you will learn.}
\end{enumerate}
\end{attention}

\newpage

\coursec{Exercices}

\courssub{Je vérifie mes connaissances}

\begin{exercice}[QCM : la bonne forme]
Entoure la bonne réponse.
\begin{enumerate}
\item This exercise is getting \ldots
\begin{itemize}
\item[a)] more difficult and more difficult
\item[b)] more and more difficult
\item[c)] difficulter and difficulter
\end{itemize}
\item My brother is \ldots as me.
\begin{itemize}
\item[a)] so tall
\item[b)] as tall
\item[c)] tall
\end{itemize}
\item The \ldots you practise, the \ldots you speak.
\begin{itemize}
\item[a)] more / better
\item[b)] most / best
\item[c)] more / more good
\end{itemize}
\item Douala is \ldots Yaoundé.
\begin{itemize}
\item[a)] twice bigger than
\item[b)] two times bigger as
\item[c)] twice as big than
\end{itemize}
\end{enumerate}
\end{exercice}

\begin{exercice}[Vrai ou faux ?]
Dis si les phrases suivantes sont correctes (\emph{true}) ou incorrectes (\emph{false}) et justifie ta réponse en citant la règle de grammaire concernée.
\begin{enumerate}
\item My sister is as tall than me.
\item More you read, more you learn.
\item The sooner, the better.
\item This bag is three times as heavy as mine.
\item Life in the city is more and more expensive.
\end{enumerate}
\end{exercice}

\begin{exercice}[Vocabulaire de la comparaison]
Complète chaque phrase avec le mot ou l'expression qui convient, choisi dans la liste : \eng{twice — equal — inferior — similar — unlike — as much}.
\begin{enumerate}
\item A teacher's salary is not \ldots to a doctor's salary.
\item My new phone is \ldots to yours, but the colour is different.
\item This shop sells \ldots many goods as the one near the market.
\item \ldots his brother, Peter prefers farming to trading.
\item A bicycle is far \ldots to a bendskin in speed.
\item Our classroom is \ldots the size of the staffroom.
\end{enumerate}
\end{exercice}

\begin{exercice}[Formes à compléter]
Mets l'adjectif entre parenthèses à la forme qui convient (comparatif, superlatif, égalité, proportion).
\begin{enumerate}
\item The \ldots (hard) you work, the \ldots (good) your results will be.
\item This exercise is \ldots (easy) than the last one.
\item My grandmother is not as \ldots (strong) as she used to be.
\item The \ldots (much) you save, the \ldots (little) you will borrow.
\item Cocoa is getting \ldots and \ldots (expensive) on the world market.
\end{enumerate}
\end{exercice}

\courssub{Je m'entraîne}

\begin{exercice}[Traduction vers l'anglais]
Traduis les phrases suivantes en anglais correct.
\begin{enumerate}
\item Le riz devient de plus en plus cher au Cameroun.
\item Ce marché est deux fois plus grand que l'autre.
\item Plus tu étudies, meilleures seront tes notes.
\item Mon père gagne moins d'argent que mon oncle.
\item Ce taxi est aussi rapide qu'un bendskin.
\end{enumerate}
\end{exercice}

\begin{exercice}[Transformation de phrases]
Réécris chaque phrase en commençant par les mots indiqués, sans changer le sens.
\begin{enumerate}
\item A doctor earns more money than a teacher. \\
A teacher does not \ldots
\item No other student in the class is as hardworking as Amina. \\
Amina is \ldots
\item The lorry is cheaper than the train. \\
The train is not \ldots
\item As the price of fuel rises, transport fares rise too. \\
The higher \ldots
\item My house is big. Your house is big too. \\
My house is as \ldots
\end{enumerate}
\end{exercice}

\begin{exercice}[Texte à trous]
Complète le texte suivant avec les mots de la liste : \eng{the more — as — twice — less — better — the most — than}.

\begin{textebox}[At the market]
Mama Béa has sold tomatoes at Mokolo market for twenty years. Business is \ldots profitable today \ldots it was before, because there are more sellers. Her stall is \ldots as big as her neighbour's, but she earns \ldots money because her prices are lower. She always says: “\ldots you smile at customers, the \ldots they buy.” Of all the traders, she is \ldots respected, because she treats everyone \ldots an equal.
\end{textebox}
\end{exercice}

\begin{exercice}[Appariements]
Associe chaque début de phrase (1--6) à la fin qui convient (a--f).
\begin{enumerate}
\item The earlier we leave home,
\item This exercise is getting more
\item A car is three times
\item My essay is not as long
\item The less you revise,
\item No occupation is
\begin{itemize}
\item[a)] as expensive as a lorry.
\item[b)] and more difficult.
\item[c)] the less traffic we meet.
\item[d)] as yours.
\item[e)] the worse your marks will be.
\item[f)] superior to another.
\end{itemize}
\end{enumerate}
\end{exercice}

\courssub{Je me prépare au Bac}

\begin{exercice}[Compréhension et langue — type Bac]
Lis le texte suivant, puis réponds aux questions.

\begin{textebox}[Two brothers, two occupations]
Ngum and Che are brothers from Bamenda. Ngum is a civil servant in Yaoundé; Che is a farmer in their village. Ngum's salary is higher than his brother's income, but his life is not easier. Everything in the capital is getting more and more expensive: rent, transport and food cost twice as much as in the village. Che earns less money, yet he is as happy as his brother, if not happier. His food is fresher, his house is bigger, and the more his farm grows, the more independent he becomes. “The richer you are in money,” Che often says, “the poorer you can be in peace of mind.” Ngum does not completely agree, but he admits that village life is far less stressful than city life.
\end{textebox}

\textbf{Comprehension questions} (réponds en anglais, par des phrases complètes) :
\begin{enumerate}
\item What are the occupations of the two brothers?
\item Give two reasons why Ngum's life in Yaoundé is not easier than Che's.
\item Why does Che become more independent as his farm grows?
\item In your own words, explain Che's saying: “The richer you are in money, the poorer you can be in peace of mind.”
\end{enumerate}

\textbf{Language work} :
\begin{enumerate}
\item Find in the text: (a) one comparison of equality; (b) one double comparative of proportion; (c) one multiplicative comparison.
\item Rewrite: “Ngum's salary is higher than his brother's income.” Begin: “Che's income is not \ldots”
\item Correct the mistake: “Che is more happier as his brother.”
\end{enumerate}
\end{exercice}

\begin{exercice}[Traduction — type Bac]
Traduis le passage suivant en anglais (environ 60 mots) :

« Au Cameroun, la vie en ville devient de plus en plus chère. Les loyers à Douala sont deux fois plus élevés que dans les petites villes. Plus les prix augmentent, plus les familles souffrent. Pourtant, un enseignant ne gagne pas autant qu'un commerçant. Beaucoup de jeunes pensent que la vie au village est moins stressante que la vie en ville. »
\end{exercice}

\begin{exercice}[Composition — type Bac]
\textbf{Topic:} \eng{Compare life in town and life in the village. Which one do you prefer and why?}

Rédige une composition d'environ \textbf{150 à 180 mots} en anglais. Tu dois :
\begin{itemize}
\item utiliser \textbf{au moins trois structures de comparaison} différentes vues dans la leçon (égalité avec \eng{as \ldots as}, proportion avec \eng{the more \ldots the more}, multiplication avec \eng{twice / three times}, gradation avec \eng{more and more}) ;
\item organiser ton texte en trois paragraphes (introduction, comparaison, opinion personnelle) ;
\item souligner les structures de comparaison utilisées.
\end{itemize}
\end{exercice}

\newpage

\coursec{Corrigés des exercices}

\courssub{Corrigés des exercices}

\begin{corrige}[Exercice 1 — QCM : la bonne forme]
\begin{enumerate}
\item \textbf{b) more and more difficult} \\
Structure de gradation croissante (double comparatif) : \eng{more and more + adjectif long}. \eng{« more difficult and more difficult »} est possible mais moins idiomatique pour exprimer une progression continue ; \eng{« difficulter »} n'existe pas.
\item \textbf{b) as tall} \\
Comparaison d'égalité : \eng{as + adjectif + as}. \eng{so tall} s'utilise dans la conséquence (\eng{so tall that...}) ou la négation (\eng{not so tall as}), pas dans l'affirmation d'égalité simple.
\item \textbf{a) more / better} \\
Proportion (corrélation) : \eng{The + comparatif ..., the + comparatif}. \eng{Good} a un comparatif irrégulier : \eng{better}. \eng{Most/best} sont des superlatifs ; \eng{more good} est incorrect.
\item \textbf{a) twice bigger than} \\
Comparaison multiplicative de supériorité : \eng{twice / three times + comparatif + than}. \eng{Option b)} : \eng{as} ne suit pas un comparatif de supériorité. \eng{Option c)} : \eng{than} ne suit pas \eng{as... as}.
\end{enumerate}
\end{corrige}

\begin{attention}[Piège classique : \eng{twice bigger than} vs \eng{twice as big as}]
En anglais standard moderne, on préfère souvent \eng{twice as big as} (égalité multiplicative) pour dire « deux fois plus grand que » (A = 2 × B). \eng{Twice bigger than} est toutefois fréquent dans les examens camerounais (et signifie logiquement « trois fois la taille », A = B + 2B), c'est souvent la réponse attendue parmi les choix proposés. Lisez bien les options : si \eng{twice as big as} est présent, c'est la forme standard à privilégier.
\end{attention}

\begin{corrige}[Exercice 2 — Vrai ou faux ?]
\begin{enumerate}
\item \textbf{False.} \eng{My sister is as tall \textbf{than} me.} $\rightarrow$ Règle : la comparaison d'égalité utilise la structure \eng{as + adjectif + \textbf{as}}.
\item \textbf{False.} \eng{\textbf{More} you read, \textbf{more} you learn.} $\rightarrow$ Règle : la proportion (corrélation) exige l'article défini \eng{\textbf{The}} devant les deux comparatifs : \eng{\textbf{The} more you read, \textbf{the} more you learn.}
\item \textbf{True.} \eng{The sooner, the better.} $\rightarrow$ Forme elliptique correcte de la proportion (verbes omis par récupération du contexte).
\item \textbf{True.} \eng{This bag is three times as heavy as mine.} $\rightarrow$ Structure multiplicative d'égalité : \eng{three times + as + adjectif + as}.
\item \textbf{True.} \eng{Life in the city is more and more expensive.} $\rightarrow$ Double comparatif d'augmentation (gradation) : \eng{more and more + adjectif long}.
\end{enumerate}
\end{corrige}

\begin{corrige}[Exercice 3 — Vocabulaire de la comparaison]
\begin{enumerate}
\item \textbf{equal} \eng{(A teacher's salary is not \textbf{equal to} a doctor's salary.)} — \eng{Equal to} = égal à.
\item \textbf{similar} \eng{(My new phone is \textbf{similar to} yours...)} — \eng{Similar to} = similaire à.
\item \textbf{as much} \eng{(This shop sells \textbf{as much} many goods as the one near the market.)} — \textbf{Analyse :} La phrase originale contient une anomalie (\eng{as much many}). \eng{Goods} est comptable pluriel, donc le quantificateur correct est \eng{many}. La structure d'égalité est \eng{as many goods as}. Si l'exercice impose de choisir \eng{as much} dans la liste pour combler un trou unique (ex. \eng{as \_\_\_\_ many goods as}), le mot de la liste est \eng{as much}, mais la phrase résultante est fautive. \textbf{Correction attendue :} L'élève écrit \eng{as much} (mot de la liste) et note que la phrase correcte est \eng{as many goods as}.
\item \textbf{Unlike} \eng{(\textbf{Unlike} his brother, Peter prefers farming to trading.)} — \eng{Unlike} = contrairement à (préposition).
\item \textbf{inferior} \eng{(A bicycle is far \textbf{inferior to} a bendskin in speed.)} — \eng{Inferior to} = inférieur à (adjectif latin, pas de \eng{than}).
\item \textbf{twice} \eng{(Our classroom is \textbf{twice} the size of the staffroom.)} — \eng{Twice the size of} = deux fois la taille de (structure multiplicative nominale).
\end{enumerate}
\end{corrige}

\begin{corrige}[Exercice 4 — Formes à compléter]
\begin{enumerate}
\item \eng{The \textbf{harder} you work, the \textbf{better} your results will be.} \\
Proportion : \eng{The + comparatif (hard $\rightarrow$ harder), the + comparatif irrégulier (good $\rightarrow$ better)}.
\item \eng{This exercise is \textbf{easier} than the last one.} \\
Comparatif de supériorité adjectif court (1 syllabe) : \eng{easy $\rightarrow$ easier} (y $\rightarrow$ i + er).
\item \eng{My grandmother is not as \textbf{strong} as she used to be.} \\
Négation de l'égalité : \eng{not as + adjectif \textbf{base} + as}.
\item \eng{The \textbf{more} you save, the \textbf{less} you will borrow.} \\
Proportion avec quantificateurs : \eng{The more (comparatif de much/many), the less (comparatif de little)}.
\item \eng{Cocoa is getting \textbf{more and more expensive} on the world market.} \\
Double comparatif d'augmentation (gradation) : \eng{more and more + adjectif long}.
\end{enumerate}
\end{corrige}

\begin{corrige}[Exercice 5 — Traduction vers l'anglais]
\begin{enumerate}
\item \eng{Rice is becoming \textbf{more and more expensive} in Cameroon.} (ou \eng{getting more and more expensive})
\item \eng{This market is \textbf{twice as big as} the other one.} (ou \eng{twice bigger than} — voir note Exercice 1.4)
\item \eng{\textbf{The more} you study, \textbf{the better} your marks will be.} (ou \eng{grades})
\item \eng{My father earns \textbf{less money than} my uncle.}
\item \eng{This taxi is \textbf{as fast as} a bendskin.}
\end{enumerate}
\end{corrige}

\begin{corrige}[Exercice 6 — Transformation de phrases]
\begin{enumerate}
\item \eng{A teacher does not earn \textbf{as much money as} a doctor.} (Égalité négative : \eng{not as much as})
\item \eng{Amina is \textbf{the most hardworking student in the class}.} (Superlatif de supériorité ; \eng{hardworking} $\rightarrow$ \eng{most hardworking})
\item \eng{The train is not \textbf{as cheap as} the lorry.} (Égalité négative ; l'antonyme de \eng{cheaper} est \eng{not as cheap as})
\item \eng{\textbf{The higher} the price of fuel \textbf{(rises)}, \textbf{the higher} transport fares \textbf{(rise)}.} (Proportion : \eng{The + comparatif ..., the + comparatif}. On peut omettre le verbe répété : \eng{The higher the price of fuel, the higher the transport fares.})
\item \eng{My house is \textbf{as big as} yours.} (Égalité affirmative)
\end{enumerate}
\end{corrige}

\begin{corrige}[Exercice 7 — Texte à trous]
\begin{enumerate}
\item \eng{less ... than} : comparatif d'infériorité (\eng{less + adj + than}).
\item \eng{twice as big as} : comparaison multiplicative d'égalité.
\item \eng{less money} : comparatif d'infériorité sur nom incomptable (\eng{less}).
\item \eng{The more ..., the more ...} : proportion (corrélation).
\item \eng{the most respected} : superlatif de supériorité (adjectif long / participe passé) avec \eng{of all}.
\item \eng{as an equal} : comparaison d'égalité (fonction) : \eng{as} = en tant que / comme.
\end{enumerate}
\end{corrige}

\begin{textebox}[At the market — Corrigé]
Mama Béa has sold tomatoes at Mokolo market for twenty years. Business is \textbf{less} profitable today \textbf{than} it was before, because there are more sellers. Her stall is \textbf{twice} as big as her neighbour's, but she earns \textbf{less} money because her prices are lower. She always says: “\textbf{The more} you smile at customers, \textbf{the more} they buy.” Of all the traders, she is \textbf{the most} respected, because she treats everyone \textbf{as} an equal.
\end{textebox}

\begin{corrige}[Exercice 8 — Appariements]
\begin{center}
\begin{tabular}{|c|c|l|}
\hline
\rowcolor{popblueL} \textbf{Début} & \textbf{Fin} & \textbf{Structure} \\
\hline
1 & c & \eng{The earlier ..., the less ...} (Proportion) \\
2 & b & \eng{more and more difficult} (Double comparatif / Gradation) \\
3 & a & \eng{three times as expensive as} (Multiplicative d'égalité) \\
4 & d & \eng{not as long as} (Égalité négative) \\
5 & e & \eng{The less ..., the worse ...} (Proportion) \\
6 & f & \eng{No ... is superior to ...} (Négation + comparatif latin) \\
\hline
\end{tabular}
\end{center}
\end{corrige}

\begin{corrige}[Exercice 9 — Compréhension et langue — type Bac]
\textbf{Barème indicatif :} 10 points (Compréhension 4 pts, Langue 6 pts).

\end{corrige}

\courssub{Comprehension questions (propositions de réponses)}
\begin{enumerate}
\item Ngum is a civil servant working in Yaoundé, while Che is a farmer in their village.
\item First, the cost of living in Yaoundé is much higher: rent, transport and food cost twice as much as in the village. Second, city life is far more stressful than village life.
\item As his farm grows, he produces more food and income, so he depends less on outside sources; he becomes more self-sufficient / independent.
\item Che means that having a lot of money does not guarantee happiness. Rich people often have many worries, responsibilities and stress, so they can lack peace of mind, unlike poorer people who may live a simpler, quieter life.
\end{enumerate}

\courssub{Language work}
\begin{enumerate}
\item \textbf{(a) Equality:} \eng{« as happy as his brother »} (l. 5). \\
\textbf{(b) Double comparative of proportion:} \eng{« the more his farm grows, the more independent he becomes »} (l. 6-7). \\
\textbf{(c) Multiplicative comparison:} \eng{« cost twice as much as »} (l. 4).
\item \eng{Che's income is not \textbf{as high as} Ngum's salary.} (ou \eng{not higher than})
\item \textbf{Mistake:} \eng{« Che is more happier as his brother. »} \\
\textbf{Correction:} \eng{Che is \textbf{happier than} his brother.} (ou \eng{as happy as}).
\end{enumerate}


\begin{attention}[Deux erreurs corrigées dans la phrase de l'élève]
\begin{enumerate}
\item \eng{More happier} : double comparatif illégal (adjectif court \eng{happy} $\rightarrow$ \eng{happier}, pas \eng{more happier}).
\item \eng{As} : après un comparatif de supériorité (\eng{happier}), on utilise \eng{than}, pas \eng{as}.
\end{enumerate}
\end{attention}

\begin{corrige}[Exercice 10 — Traduction — type Bac]
\textbf{Proposition de traduction (environ 60 mots) :}
\end{corrige}

\begin{textebox}[Traduction modèle]
In Cameroon, life in town is becoming \textbf{more and more expensive}. Rents in Douala are \textbf{twice as high as} in small towns. \textbf{The more} prices rise, \textbf{the more} families suffer. Yet, a teacher does not earn \textbf{as much as} a trader. Many young people think that life in the village is \textbf{less stressful than} life in town.
\end{textebox}

\textbf{Points de vigilance (Barème) :}
\begin{itemize}
\item \eng{More and more expensive} (gradation) — 1 pt
\item \eng{Twice as high as} (multiplicative) — 1 pt
\item \eng{The more ..., the more ...} (proportion) — 2 pts
\item \eng{As much as} (égalité quantité) — 1 pt
\item \eng{Less ... than} (infériorité) — 1 pt
\item Vocabulaire précis (\eng{rents, trader, stressful}) — 2 pts
\item Orthographe, ponctuation, fluidité — 2 pts
\end{itemize}
Total : 10 pts.

\begin{corrige}[Exercice 11 — Composition — type Bac]
\textbf{Barème indicatif :} 10 points (Contenu/Organisation 4 pts, Structures de comparaison obligatoires 3 pts, Vocabulaire/Grammaire 2 pts, Cohérence/Style 1 pt).

\textbf{Composition modèle (environ 165 mots) :}
\end{corrige}

\begin{textebox}[Town vs Village Life]
Living in town or in the village offers two very different experiences. In big cities like Douala or Yaoundé, life is \textbf{more and more expensive} and stressful. Rents are \textbf{twice as high as} in rural areas, and traffic jams make daily commutes exhausting. However, towns provide better schools, hospitals and job opportunities.

On the other hand, village life is \textbf{as peaceful as} it is affordable. Housing costs \textbf{less than} in town, and food is fresher because it comes straight from the farm. \textbf{The more} land you cultivate, \textbf{the more} independent you become. Social ties are stronger; neighbours help each other like a real family.

Personally, I \textbf{prefer village life} because I value peace of mind over material comfort. Even if I earn \textbf{less money than} a city employee, I enjoy a quality of life that money cannot buy. As the saying goes, “The richer you are in money, the poorer you can be in peace of mind.” For me, the village remains the best place to raise a family.
\end{textebox}

\textbf{Structures soulignées (vérification) :}
\begin{enumerate}
\item \eng{more and more expensive} (Gradation / Double comparatif d'augmentation)
\item \eng{twice as high as} (Comparaison multiplicative d'égalité)
\item \eng{as peaceful as} (Comparaison d'égalité)
\item \eng{less than} (Comparaison d'infériorité)
\item \eng{The more ..., the more ...} (Proportion / Double comparatif de corrélation)
\item \eng{less money than} (Comparaison d'infériorité sur nom)
\end{enumerate}
\textit{(Le texte en contient six, bien au-dessus du minimum de trois requis).}

\newpage

\coursec{Fiche bilan}

\begin{popfigure}
\begin{tikzpicture}[mindmap, grow cyclic, every node/.style={concept, align=center, text=white, font=\small}, root concept/.append style={font=\small\bfseries, fill=popdark, text width=3.2cm}, level 1/.append style={level distance=3.4cm, sibling angle=60}, level 2/.append style={level distance=2.4cm, font=\scriptsize, text width=2.2cm}]
\node[root concept]{Other forms of comparing}
  child[concept color=popblue]{ node {Equality / Inequality}
    child { node[fill=popblue!70]{as \ldots\ as / not so \ldots\ as} }
    child { node[fill=popblue!70]{the same \ldots\ as} }
  }
  child[concept color=poporange]{ node {Degree}
    child { node[fill=poporange!70]{much, far, a lot + comp.} }
    child { node[fill=poporange!70]{slightly, a bit + comp.} }
  }
  child[concept color=popgreen]{ node {Proportion}
    child { node[fill=popgreen!70]{the more \ldots\ the more} }
  }
  child[concept color=poppurple]{ node {Double comparison}
    child { node[fill=poppurple!70]{more and more / better and better} }
  }
  child[concept color=poppink]{ node {Multiplication}
    child { node[fill=poppink!70]{twice / three times as \ldots\ as} }
  }
  child[concept color=popgold]{ node {Traps}
    child { node[fill=popgold!70]{no ``more better''} }
  };
\end{tikzpicture}
\legende{Carte mentale : les autres formes de comparaison en anglais}
\end{popfigure}

\begin{bilan}
\textbf{1. Lexique clé (anglais / français)}
\begin{itemize}
\item \eng{as \ldots\ as} : aussi \ldots\ que ; \eng{not as/so \ldots\ as} : pas aussi \ldots\ que
\item \eng{the same \ldots\ as} : le même \ldots\ que ; \eng{different from} : différent de
\item \eng{twice / three times as much (many) \ldots\ as} : deux / trois fois plus \ldots\ que
\item \eng{the more \ldots\ the more} : plus \ldots\ plus ; \eng{the less \ldots\ the less} : moins \ldots\ moins
\item \eng{more and more} : de plus en plus ; \eng{less and less} : de moins en moins
\item \eng{by far the best} : de loin le meilleur ; \eng{much cheaper} : bien moins cher
\end{itemize}

\textbf{2. Règles à connaître par cœur}
\begin{center}
\begin{tabularx}{\linewidth}{|l|X|X|}
\hline
\rowcolor{popblueL} \textbf{Structure} & \textbf{Forme} & \textbf{Exemple} \\
\hline
Égalité & as + adj./adv. + as & \eng{This cocoa is as good as that one.} \\
\hline
Inégalité & not as/so + adj. + as & \eng{Buea is not as hot as Douala.} \\
\hline
Degré & much / far / a lot / slightly / a bit + comparatif & \eng{Petrol is much more expensive now.} \\
\hline
Proportion & the + comparatif + S + V, the + comparatif + S + V & \eng{The harder you work, the more you earn.} \\
\hline
Double comparatif & comparatif + and + comparatif & \eng{Life is getting more and more expensive.} \\
\hline
Multiplication & twice / three times + as \ldots\ as (ou + comparatif + than) & \eng{Yaoundé is twice as big as Bafoussam.} \\
\hline
Identité & the same + nom + as & \eng{She earns the same salary as her colleague.} \\
\hline
\end{tabularx}
\end{center}

\textbf{3. Méthodes express}
\begin{itemize}
\item Pour exprimer l'\cle{égalité} : encadrer l'adjectif court avec \eng{as \ldots\ as} ; nier avec \eng{not as/so \ldots\ as}.
\item Pour exprimer la \cle{proportion} : deux propositions, chacune introduite par \eng{the} + comparatif ; jamais de virgule oubliée entre les deux.
\item Pour \cle{multiplier} : \eng{times} prend toujours un \textbf{-s} (sauf \eng{once}, \eng{twice}) ; on peut dire \eng{three times bigger than} ou \eng{three times as big as}.
\item Pour nuancer un comparatif : \eng{much/far/a lot} (grande différence), \eng{a bit/slightly} (petite différence) — jamais \eng{very} devant un comparatif.
\end{itemize}
\end{bilan}

\begin{aretenir}[Checklist avant le Bac]
\begin{itemize}
\item[$\square$] Je sais construire la comparaison d'égalité : \eng{as + adjectif + as}.
\item[$\square$] Je sais nier l'égalité avec \eng{not as \ldots\ as} ou \eng{not so \ldots\ as}.
\item[$\square$] Je sais dire « différent de » : \eng{different from} (jamais \eng{different than} en anglais standard).
\item[$\square$] Je sais nuancer un comparatif avec \eng{much, far, a lot, slightly, a bit}.
\item[$\square$] Je n'écris jamais \eng{very better} ni \eng{more better} : un seul marqueur de comparaison.
\item[$\square$] Je maîtrise la proportion : \eng{The more you save, the richer you become.}
\item[$\square$] Je sais exprimer « de plus en plus » : \eng{more and more} / \eng{better and better}.
\item[$\square$] Je sais multiplier : \eng{twice as expensive as}, \eng{three times the price}.
\item[$\square$] Je sais dire « le même \ldots\ que » : \eng{the same \ldots\ as}.
\item[$\square$] Je connais les comparatifs irréguliers : \eng{good $\rightarrow$ better}, \eng{bad $\rightarrow$ worse}, \eng{far $\rightarrow$ further/farther}.
\item[$\square$] Je place correctement \eng{by far} devant le superlatif : \eng{by far the richest trader in the market}.
\item[$\square$] Je peux utiliser ces structures dans un paragraphe sur la vie économique (prix, salaires, métiers).
\end{itemize}
\end{aretenir}

\findecours

\end{document}
